\documentclass{article}
\usepackage{amsmath,amssymb,amsthm}
\IfFileExists{hyperref.sty}{%
  \IfFileExists{color.sty}{\usepackage{color}\definecolor{paperlinkblue}{rgb}{0.10,0.25,0.42}%
    \usepackage[colorlinks=true,linkcolor=paperlinkblue]{hyperref}}{\usepackage[hidelinks]{hyperref}}%
}{}
\ifdefined\hyperlink
  \newcommand{\paperanchor}[1]{\hypertarget{#1}{}}%
  \newcommand{\paperlink}[2]{\hyperlink{#1}{\begin{NoHyper}#2\end{NoHyper}}}%
\else
  \ifdefined\pdfextension
    \ifnum\outputmode>0
      \def\paperdest{\pdfextension dest}\def\paperstart{\pdfextension startlink}%
      \def\paperend{\pdfextension endlink}\def\papercolor{\pdfextension colorstack}%
    \fi
  \else\ifdefined\pdfstartlink
    \ifnum\pdfoutput>0
      \let\paperdest\pdfdest\let\paperstart\pdfstartlink\let\paperend\pdfendlink\let\papercolor\pdfcolorstack
    \fi
  \fi\fi
  \ifdefined\paperdest
    \newcommand{\paperanchor}[1]{\paperdest name{#1} xyz}%
    \newcommand{\paperlink}[2]{\leavevmode\paperstart attr{/Border [0 0 0]} goto name{#1}%
      \papercolor 0 push {0.10 0.25 0.42 rg 0.10 0.25 0.42 RG}#2\papercolor 0 pop\paperend}%
  \else\ifdefined\XeTeXversion
    \newcommand{\paperanchor}[1]{\special{pdf:dest (#1) [@thispage /XYZ null @ypos null]}}%
    \newcommand{\paperlink}[2]{\leavevmode\special{pdf:bann << /Subtype /Link /Border [0 0 0] /A << /S /GoTo /D (#1) >> >>}%
      \special{color push rgb 0.10 0.25 0.42}#2\special{color pop}\special{pdf:eann}}%
  \else
    \PackageError{lean-to-latex}{No PDF hyperlink backend available}{Use XeLaTeX, pdfLaTeX, LuaLaTeX, or install hyperref.}%
  \fi\fi
\fi
\newcommand{\paperlineanchor}[1]{\raisebox{10pt}[0pt][0pt]{\paperanchor{#1}}}
\DeclareRobustCommand{\paperresultname}[1]{\paperlineanchor{theorem:\arabic{theorem}}#1}
\newtheoremstyle{paperresult}{8pt}{5pt}{\normalfont\raggedright}{}{\bfseries}{.}{\newline}{}
\theoremstyle{paperresult}
\newtheorem{theorem}{\paperresultname{Theorem}}
\newcounter{paperpremise}
\newtheorem{lemma}[theorem]{\paperresultname{Lemma}}
\newtheorem{corollary}[theorem]{\paperresultname{Corollary}}
\newcounter{definingequation}
\newenvironment{proofstep}{\par\smallskip\ifdim\leftskip<3em\relax\advance\leftskip by 1em\relax\fi\noindent\ignorespaces}{\par\smallskip}
\setlength{\textwidth}{6.5in}
\setlength{\oddsidemargin}{0in}
\AddToHook{env/proof/begin}{\raggedright}
\setlength{\emergencystretch}{\linewidth}
\title{The inclusion-exclusion principle}
\author{}
\date{}
\begin{document}
\maketitle
\subsection*{Prerequisites}
\begin{description}
\item[\paperlineanchor{premise:436c6173736963616c2e6e6f745f666f72616c6c}\paperlink{premise:436c6173736963616c2e6e6f745f666f72616c6c}{P1.}] \refstepcounter{paperpremise}\label{premise:436c6173736963616c2e6e6f745f666f72616c6c}\textbf{Not forall.} Let \(\alpha\) be a sort. Let \(p\) be a function of type \(\alpha \to \mathsf{Prop}\). Then \(\neg (\forall x\, :\, \alpha,\; p(x)) \iff (\exists x\, :\, \alpha,\; \neg p(x))\).
\item[\paperlineanchor{premise:46696e7365742e636f655f656d707479}\paperlink{premise:46696e7365742e636f655f656d707479}{P2.}] \refstepcounter{paperpremise}\label{premise:46696e7365742e636f655f656d707479}\textbf{Coercion empty.} Let \(\alpha\) be a type. Then \(\operatorname{underlying\,set}(\varnothing) = \varnothing\).
\item[\paperlineanchor{premise:46696e7365742e636f6d706c5f737570}\paperlink{premise:46696e7365742e636f6d706c5f737570}{P3.}] \refstepcounter{paperpremise}\label{premise:46696e7365742e636f6d706c5f737570}\textbf{Compl sup.} Let \(\alpha\) and \(\iota\) be types. Fix a structure \(\mathit{inst}\) of type \(\operatorname{Boolean\,Algebra}\left(\alpha\right)\). Let \(s\) be a finite subset of \(\iota\). Let \(f\) be a function of type \(\iota \to \alpha\). Then \(\operatorname{compl}(\operatorname{sup}(s,\allowbreak{} f)) = \operatorname{inf}(s,\allowbreak{} i \mapsto \operatorname{compl}(f(i)))\).
\item[\paperlineanchor{premise:46696e7365742e656d7074795f737562736574}\paperlink{premise:46696e7365742e656d7074795f737562736574}{P4.}] \refstepcounter{paperpremise}\label{premise:46696e7365742e656d7074795f737562736574}\textbf{Empty subset.} Let \(\alpha\) be a type. Let \(s\) be a finite subset of \(\alpha\). Then \(\varnothing \subseteq s\).
\item[\paperlineanchor{premise:46696e7365742e65715f656d7074795f6f725f6e6f6e656d707479}\paperlink{premise:46696e7365742e65715f656d7074795f6f725f6e6f6e656d707479}{P5.}] \refstepcounter{paperpremise}\label{premise:46696e7365742e65715f656d7074795f6f725f6e6f6e656d707479}\textbf{Equals empty or nonempty.} Let \(\alpha\) be a type. Let \(s\) be a finite subset of \(\alpha\). Then \(s = \varnothing \lor s\ne\varnothing\).
\item[\paperlineanchor{premise:46696e7365742e657874}\paperlink{premise:46696e7365742e657874}{P6.}] \refstepcounter{paperpremise}\label{premise:46696e7365742e657874}\textbf{Ext.} Let \(\alpha\) be a type. Let \(s_{1}\) and \(s_{2}\) be finite subsets of \(\alpha\). Assume \(\forall a\, :\, \alpha,\; a \in s_{1} \iff a \in s_{2}\). Then \(s_{1} = s_{2}\).
\item[\paperlineanchor{premise:46696e7365742e66696c7465725f657127}\paperlink{premise:46696e7365742e66696c7465725f657127}{P7.}] \refstepcounter{paperpremise}\label{premise:46696e7365742e66696c7465725f657127}\textbf{Filter eq'.} Let \(\beta\) be a type. Assume equality on \(\beta\) is decidable. Let \(s\) be a finite subset of \(\beta\). Let \(b\) be an element of \(\beta\). Then \(\{a\in s \mid a = b\} = \begin{cases}\{b\} &\text{if } b \in s\\ \varnothing &\text{otherwise}\end{cases}\).
\item[\paperlineanchor{premise:46696e7365742e696e66275f65715f696e66}\paperlink{premise:46696e7365742e696e66275f65715f696e66}{P8.}] \refstepcounter{paperpremise}\label{premise:46696e7365742e696e66275f65715f696e66}\textbf{Inf' equals inf.} Let \(\alpha\) and \(\beta\) be types. Fix a structure \(\mathit{inst}\) of type \(\operatorname{Semilattice\,Inf}\left(\alpha\right)\). Fix a structure \(\mathit{inst}_{1}\) of type \(\operatorname{Order\,Top}\left(\alpha\right)\). Let \(s\) be a finite subset of \(\beta\). Let \(H\) be a proof of \(s\ne\varnothing\). Let \(f\) be a function of type \(\beta \to \alpha\). Then \(\operatorname{inf^{\prime}}\left(s, f\right) = \operatorname{inf}\left(s, f\right)\).
\item[\paperlineanchor{premise:46696e7365742e6d656d5f6269556e696f6e}\paperlink{premise:46696e7365742e6d656d5f6269556e696f6e}{P9.}] \refstepcounter{paperpremise}\label{premise:46696e7365742e6d656d5f6269556e696f6e}\textbf{Membership bi union.} Let \(\alpha\) and \(\beta\) be types. Let \(s\) be a finite subset of \(\alpha\). Let \(t\) be a family of finite subsets of \(\beta\) indexed by \(\alpha\). Assume equality on \(\beta\) is decidable. Let \(b\) be an element of \(\beta\). Then \(b \in \bigcup_{i\in s} t(i) \iff \left(\exists a\, :\, \alpha,\; a \in s \land b \in t\left(a\right)\right)\).
\item[\paperlineanchor{premise:46696e7365742e6d656d5f636f65}\paperlink{premise:46696e7365742e6d656d5f636f65}{P10.}] \refstepcounter{paperpremise}\label{premise:46696e7365742e6d656d5f636f65}\textbf{Membership coercion.} Let \(\alpha\) be a type. Let \(a\) be an element of \(\alpha\). Let \(s\) be a finite subset of \(\alpha\). Then \(a \in \operatorname{underlying\,set}(s) \iff a \in s\).
\item[\paperlineanchor{premise:46696e7365742e6d656d5f66696c746572}\paperlink{premise:46696e7365742e6d656d5f66696c746572}{P11.}] \refstepcounter{paperpremise}\label{premise:46696e7365742e6d656d5f66696c746572}\textbf{Membership filter.} Let \(\alpha\) be a type. Let \(p\) be a function of type \(\alpha \to \mathsf{Prop}\). Assume the predicate \(p\) is decidable. Let \(s\) be a finite subset of \(\alpha\). Let \(a\) be an element of \(\alpha\). Then \(a \in \{x\in s \mid p(x)\} \iff a \in s \land p\left(a\right)\).
\item[\paperlineanchor{premise:46696e7365742e6d656d5f696e6627}\paperlink{premise:46696e7365742e6d656d5f696e6627}{P12.}] \refstepcounter{paperpremise}\label{premise:46696e7365742e6d656d5f696e6627}\textbf{Membership inf'.} Let \(\alpha\) and \(\iota\) be types. Assume equality on \(\alpha\) is decidable. Let \(s\) be a finite subset of \(\iota\). Let \(f\) be a family of finite subsets of \(\alpha\) indexed by \(\iota\). Let \(a\) be an element of \(\alpha\). Let \(\mathit{hs}\) be a proof of \(s\ne\varnothing\). Then \(a \in \bigcap_{i\in s} f(i) \iff \left(\forall i\, :\, \iota,\; i \in s \implies a \in f\left(i\right)\right)\).
\item[\paperlineanchor{premise:46696e7365742e6d656d5f706f776572736574}\paperlink{premise:46696e7365742e6d656d5f706f776572736574}{P13.}] \refstepcounter{paperpremise}\label{premise:46696e7365742e6d656d5f706f776572736574}\textbf{Membership power set.} Let \(\alpha\) be a type. Let \(s\) and \(t\) be finite subsets of \(\alpha\). Then \(s \in \mathcal P(t) \iff s \subseteq t\).
\item[\paperlineanchor{premise:46696e7365742e6e6f745f6e6f6e656d7074795f6966665f65715f656d707479}\paperlink{premise:46696e7365742e6e6f745f6e6f6e656d7074795f6966665f65715f656d707479}{P14.}] \refstepcounter{paperpremise}\label{premise:46696e7365742e6e6f745f6e6f6e656d7074795f6966665f65715f656d707479}\textbf{Not nonempty if and only if equals empty.} Let \(\alpha\) be a type. Let \(s\) be a finite subset of \(\alpha\). Then \(\neg \left(s\ne\varnothing\right) \iff s = \varnothing\).
\item[\paperlineanchor{premise:46696e7365742e70726f645f636f6e6772}\paperlink{premise:46696e7365742e70726f645f636f6e6772}{P15.}] \refstepcounter{paperpremise}\label{premise:46696e7365742e70726f645f636f6e6772}\textbf{Equality of products with equal factors.} Let \(\iota\) and \(M\) be types. Let \(s_{1}\) and \(s_{2}\) be finite subsets of \(\iota\). Fix a structure \(\mathit{inst}\) of type \(\operatorname{Comm\,Monoid}\left(M\right)\). Let \(f\) and \(g\) be functions of type \(\iota \to M\). Assume \(s_{1} = s_{2}\). Assume \(\forall x\, :\, \iota,\; x \in s_{2} \implies f\left(x\right) = g\left(x\right)\). Then \(\prod_{i\in s_{1}}\left(f(i)\right) = \prod_{i\in s_{2}}\left(g(i)\right)\).
\item[\paperlineanchor{premise:46696e7365742e70726f645f636f6e73745f6f6e65}\paperlink{premise:46696e7365742e70726f645f636f6e73745f6f6e65}{P16.}] \refstepcounter{paperpremise}\label{premise:46696e7365742e70726f645f636f6e73745f6f6e65}\textbf{Product constant one.} Let \(\alpha\) and \(\beta\) be types. Let \(s\) be a finite subset of \(\alpha\). Fix a structure \(\mathit{inst}\) of type \(\operatorname{Comm\,Monoid}\left(\beta\right)\). Then \(\prod_{x\in s}\left(1\right) = 1\).
\item[\paperlineanchor{premise:46696e7365742e70726f645f65715f7a65726f}\paperlink{premise:46696e7365742e70726f645f65715f7a65726f}{P17.}] \refstepcounter{paperpremise}\label{premise:46696e7365742e70726f645f65715f7a65726f}\textbf{Product equals zero.} Let \(\iota\) and \(M_{0}\) be types. Fix a structure \(\mathit{inst}\) of type \(\operatorname{Comm\,Monoid\,With\,Zero}\left(M_{0}\right)\). Let \(f\) be a function of type \(\iota \to M_{0}\). Let \(s\) be a finite subset of \(\iota\). Let \(i\) be an element of \(\iota\). Assume \(i \in s\). Assume \(f\left(i\right) = 0\). Then \(\prod_{j\in s}\left(f\left(j\right)\right) = 0\).
\item[\paperlineanchor{premise:46696e7365742e70726f645f737562}\paperlink{premise:46696e7365742e70726f645f737562}{P18.}] \refstepcounter{paperpremise}\label{premise:46696e7365742e70726f645f737562}\textbf{Product difference.} Let \(\iota\) and \(\alpha\) be types. Fix a structure \(\mathit{inst}\) of type \(\operatorname{Comm\,Ring}\left(\alpha\right)\). Assume equality on \(\iota\) is decidable. Let \(f\) and \(g\) be functions of type \(\iota \to \alpha\). Let \(s\) be a finite subset of \(\iota\). Then \(\prod_{i\in s}(f(i) - g(i)) = \sum_{t\in \mathcal P(s)}((-1)^{|t|} \cdot \prod_{i\in s \setminus t}(f(i)) \cdot \prod_{i\in t}(g(i)))\).
\item[\paperlineanchor{premise:46696e7365742e73646966665f656d707479}\paperlink{premise:46696e7365742e73646966665f656d707479}{P19.}] \refstepcounter{paperpremise}\label{premise:46696e7365742e73646966665f656d707479}\textbf{Sdiff empty.} Let \(\alpha\) be a type. Assume equality on \(\alpha\) is decidable. Let \(s\) be a finite subset of \(\alpha\). Then \(s \setminus \varnothing = s\).
\item[\paperlineanchor{premise:46696e7365742e736d756c5f73756d}\paperlink{premise:46696e7365742e736d756c5f73756d}{P20.}] \refstepcounter{paperpremise}\label{premise:46696e7365742e736d756c5f73756d}\textbf{Smul sum.} Let \(\alpha\), \(\beta\), and \(\gamma\) be types. Equip \(\beta\) with a commutative additive monoid structure. Fix a structure \(\mathit{inst}_{1}\) of type \(\operatorname{Distrib\,SMul}\left(\alpha, \beta\right)\). Let \(r\) be an element of \(\alpha\). Let \(f\) be a function of type \(\gamma \to \beta\). Let \(s\) be a finite subset of \(\gamma\). Then \(r \mathbin{\bullet} \sum_{x\in s}(f(x)) = \sum_{x\in s}(r \mathbin{\bullet} f(x))\).
\item[\paperlineanchor{premise:46696e7365742e7375627365745f6269556e696f6e5f6f665f6d656d}\paperlink{premise:46696e7365742e7375627365745f6269556e696f6e5f6f665f6d656d}{P21.}] \refstepcounter{paperpremise}\label{premise:46696e7365742e7375627365745f6269556e696f6e5f6f665f6d656d}\textbf{Subset bi union of membership.} Let \(\alpha\) and \(\beta\) be types. Let \(s\) be a finite subset of \(\alpha\). Assume equality on \(\beta\) is decidable. Let \(u\) be a family of finite subsets of \(\beta\) indexed by \(\alpha\). Let \(x\) be an element of \(\alpha\). Assume \(x \in s\). Then \(u\left(x\right) \subseteq \bigcup_{i\in s} u(i)\).
\item[\paperlineanchor{premise:46696e7365742e73756d5f617474616368}\paperlink{premise:46696e7365742e73756d5f617474616368}{P22.}] \refstepcounter{paperpremise}\label{premise:46696e7365742e73756d5f617474616368}\textbf{Sum attach.} Let \(\iota\) and \(M\) be types. Equip \(M\) with a commutative additive monoid structure. Let \(s\) be a finite subset of \(\iota\). Let \(f\) be a function of type \(\iota \to M\). Then \(\sum_{x\in \operatorname{attach}(s)}(f(x)) = \sum_{x\in s}(f(x))\).
\item[\paperlineanchor{premise:46696e7365742e73756d5f636f6d6d}\paperlink{premise:46696e7365742e73756d5f636f6d6d}{P23.}] \refstepcounter{paperpremise}\label{premise:46696e7365742e73756d5f636f6d6d}\textbf{Interchanging two finite sums.} Let \(\alpha\), \(\beta\), and \(\gamma\) be types. Equip \(\beta\) with a commutative additive monoid structure. Let \(s\) be a finite subset of \(\gamma\). Let \(t\) be a finite subset of \(\alpha\). Let \(f\) be a function of type \(\gamma \to \alpha \to \beta\). Then \(\sum_{x\in s}(\sum_{y\in t}(f(x,\allowbreak{} y))) = \sum_{y\in t}(\sum_{x\in s}(f(x,\allowbreak{} y)))\).
\item[\paperlineanchor{premise:46696e7365742e73756d5f636f6d706c5f6164645f73756d}\paperlink{premise:46696e7365742e73756d5f636f6d706c5f6164645f73756d}{P24.}] \refstepcounter{paperpremise}\label{premise:46696e7365742e73756d5f636f6d706c5f6164645f73756d}\textbf{Sum compl sum sum.} Let \(\iota\) and \(M\) be types. Equip \(M\) with a commutative additive monoid structure. Fix a finite enumeration of \(\iota\). Assume equality on \(\iota\) is decidable. Let \(s\) be a finite subset of \(\iota\). Let \(f\) be a function of type \(\iota \to M\). Then \(\sum_{i\in s^{\mathrm c}}(f(i)) + \sum_{i\in s}(f(i)) = \sum_{i\in \iota}(f(i))\).
\item[\paperlineanchor{premise:46696e7365742e73756d5f636f6e6772}\paperlink{premise:46696e7365742e73756d5f636f6e6772}{P25.}] \refstepcounter{paperpremise}\label{premise:46696e7365742e73756d5f636f6e6772}\textbf{Equality of sums with equal summands.} Let \(\iota\) and \(M\) be types. Let \(s_{1}\) and \(s_{2}\) be finite subsets of \(\iota\). Equip \(M\) with a commutative additive monoid structure. Let \(f\) and \(g\) be functions of type \(\iota \to M\). Assume \(s_{1} = s_{2}\). Assume \(\forall x\, :\, \iota,\; x \in s_{2} \implies f\left(x\right) = g\left(x\right)\). Then \(\sum_{i\in s_{1}}\left(f(i)\right) = \sum_{i\in s_{2}}\left(g(i)\right)\).
\item[\paperlineanchor{premise:46696e7365742e73756d5f636f6e7374}\paperlink{premise:46696e7365742e73756d5f636f6e7374}{P26.}] \refstepcounter{paperpremise}\label{premise:46696e7365742e73756d5f636f6e7374}\textbf{The sum of a constant.} Let \(\iota\) and \(M\) be types. Let \(s\) be a finite subset of \(\iota\). Equip \(M\) with a commutative additive monoid structure. Let \(b\) be an element of \(M\). Then \(\sum_{x\in s}\left(b\right) = \sum_{k=1}^{\left|s\right|}\left(b\right)\).
\item[\paperlineanchor{premise:46696e7365742e73756d5f636f6e73745f7a65726f}\paperlink{premise:46696e7365742e73756d5f636f6e73745f7a65726f}{P27.}] \refstepcounter{paperpremise}\label{premise:46696e7365742e73756d5f636f6e73745f7a65726f}\textbf{Sum constant zero.} Let \(\alpha\) and \(\beta\) be types. Let \(s\) be a finite subset of \(\alpha\). Equip \(\beta\) with a commutative additive monoid structure. Then \(\sum_{x\in s}\left(0\right) = 0\).
\item[\paperlineanchor{premise:46696e7365742e73756d5f64697465}\paperlink{premise:46696e7365742e73756d5f64697465}{P28.}] \refstepcounter{paperpremise}\label{premise:46696e7365742e73756d5f64697465}\textbf{Sum dite.} Let \(\alpha\) and \(\beta\) be types. Equip \(\beta\) with a commutative additive monoid structure. Let \(s\) be a finite subset of \(\alpha\). Let \(p\) be a function of type \(\alpha \to \mathsf{Prop}\). Let \(\mathit{hp}\) be an element of \(\operatorname{Decidable\,Pred}\left(p\right)\). Let \(f\) be a function of type \(\Pi_{x\, :\, \alpha}\, p\left(x\right) \to \beta\). Let \(g\) be a function of type \(\Pi_{x\, :\, \alpha}\, \neg p\left(x\right) \to \beta\). Then \(\sum_{x\in s}(\begin{cases}f(x,\allowbreak{} \mathit{hx}) &\text{given } \mathit{hx}\text{ proving } p(x)\\ g(x,\allowbreak{} \mathit{hx}) &\text{given } \mathit{hx}\text{ proving } \neg p(x)\end{cases}) = \sum_{x\in \{x\in s \mid p(x)\}}(f(x,\allowbreak{} \operatorname{right}(\operatorname{Iff\mathord{:}mp}(\operatorname{mem\,filter},\allowbreak{} \operatorname{property}(x))))) + \sum_{x\in \{x\in s \mid \neg p(x)\}}(g(x,\allowbreak{} \operatorname{right}(\operatorname{Iff\mathord{:}mp}(\operatorname{mem\,filter},\allowbreak{} \operatorname{property}(x)))))\).
\item[\paperlineanchor{premise:46696e7365742e73756d5f66696c7465725f6e6f745f6164645f73756d5f66696c746572}\paperlink{premise:46696e7365742e73756d5f66696c7465725f6e6f745f6164645f73756d5f66696c746572}{P29.}] \refstepcounter{paperpremise}\label{premise:46696e7365742e73756d5f66696c7465725f6e6f745f6164645f73756d5f66696c746572}\textbf{Sum filter not sum sum filter.} Let \(\iota\) and \(M\) be types. Equip \(M\) with a commutative additive monoid structure. Let \(s\) be a finite subset of \(\iota\). Let \(p\) be a function of type \(\iota \to \mathsf{Prop}\). Assume the predicate \(p\) is decidable. Fix a structure \(\mathit{inst}_{2}\) of type \(\Pi_{x\, :\, \iota}\, \operatorname{Decidable}\left(\neg p\left(x\right)\right)\). Let \(f\) be a function of type \(\iota \to M\). Then \(\sum_{x\in \{x\in s \mid \neg p(x)\}}(f(x)) + \sum_{x\in \{x\in s \mid p(x)\}}(f(x)) = \sum_{x\in s}(f(x))\).
\item[\paperlineanchor{premise:46696e7365742e73756d5f697465}\paperlink{premise:46696e7365742e73756d5f697465}{P30.}] \refstepcounter{paperpremise}\label{premise:46696e7365742e73756d5f697465}\textbf{Sum ite.} Let \(\alpha\) and \(\beta\) be types. Equip \(\beta\) with a commutative additive monoid structure. Let \(s\) be a finite subset of \(\alpha\). Let \(p\) be a function of type \(\alpha \to \mathsf{Prop}\). Let \(\mathit{hp}\) be an element of \(\operatorname{Decidable\,Pred}\left(p\right)\). Let \(f\) and \(g\) be functions of type \(\alpha \to \beta\). Then \(\sum_{x\in s}(\begin{cases}f(x) &\text{if } p(x)\\ g(x) &\text{otherwise}\end{cases}) = \sum_{x\in \{x\in s \mid p(x)\}}(f(x)) + \sum_{x\in \{x\in s \mid \neg p(x)\}}(g(x))\).
\item[\paperlineanchor{premise:46696e7365742e73756d5f6e65675f64697374726962}\paperlink{premise:46696e7365742e73756d5f6e65675f64697374726962}{P31.}] \refstepcounter{paperpremise}\label{premise:46696e7365742e73756d5f6e65675f64697374726962}\textbf{Sum neg distrib.} Let \(\alpha\) and \(\beta\) be types. Let \(s\) be a finite subset of \(\alpha\). Let \(f\) be a function of type \(\alpha \to \beta\). Fix a structure \(\mathit{inst}\) of type \(\operatorname{Subtraction\,Comm\,Monoid}\left(\beta\right)\). Then \(\sum_{x\in s}\left(-f\left(x\right)\right) = -\left(\sum_{x\in s}\left(f\left(x\right)\right)\right)\).
\item[\paperlineanchor{premise:46696e7365742e73756d5f73696e676c65746f6e}\paperlink{premise:46696e7365742e73756d5f73696e676c65746f6e}{P32.}] \refstepcounter{paperpremise}\label{premise:46696e7365742e73756d5f73696e676c65746f6e}\textbf{Sum singleton.} Let \(\iota\) and \(M\) be types. Equip \(M\) with a commutative additive monoid structure. Let \(f\) be a function of type \(\iota \to M\). Let \(a\) be an element of \(\iota\). Then \(\sum_{x\in \{a\}}\left(f\left(x\right)\right) = f\left(a\right)\).
\item[\paperlineanchor{premise:46696e7365742e73756d5f736d756c}\paperlink{premise:46696e7365742e73756d5f736d756c}{P33.}] \refstepcounter{paperpremise}\label{premise:46696e7365742e73756d5f736d756c}\textbf{Sum smul.} Let \(\iota\), \(R\), and \(M\) be types. Equip \(R\) with a semiring structure. Equip \(M\) with a commutative additive monoid structure. Fix a structure \(\mathit{inst}_{2}\) of type \(\operatorname{Module}\left(R, M\right)\). Let \(f\) be a function of type \(\iota \to R\). Let \(s\) be a finite subset of \(\iota\). Let \(x\) be an element of \(M\). Then \(\sum_{i\in s}(f(i)) \mathbin{\bullet} x = \sum_{i\in s}(f(i) \mathbin{\bullet} x)\).
\item[\paperlineanchor{premise:46696e7365742e7375705f65715f6269556e696f6e}\paperlink{premise:46696e7365742e7375705f65715f6269556e696f6e}{P34.}] \refstepcounter{paperpremise}\label{premise:46696e7365742e7375705f65715f6269556e696f6e}\textbf{Sup equals bi union.} Let \(\alpha\) and \(\beta\) be types. Assume equality on \(\beta\) is decidable. Let \(s\) be a finite subset of \(\alpha\). Let \(t\) be a family of finite subsets of \(\beta\) indexed by \(\alpha\). Then \(\operatorname{sup}\left(s, t\right) = \bigcup_{i\in s} t(i)\).
\item[\paperlineanchor{premise:4d756c5a65726f436c6173732e6d756c5f7a65726f}\paperlink{premise:4d756c5a65726f436c6173732e6d756c5f7a65726f}{P35.}] \refstepcounter{paperpremise}\label{premise:4d756c5a65726f436c6173732e6d756c5f7a65726f}\textbf{Product zero.} Let \(M_{0}\) be a type. Fix a structure \(\mathit{self}\) of type \(\operatorname{Mul\,Zero\,Class}\left(M_{0}\right)\). Let \(a\) be an element of \(M_{0}\). Then \(a \cdot 0 = 0\).
\item[\paperlineanchor{premise:50692e696e737453756273696e676c65746f6e}\paperlink{premise:50692e696e737453756273696e676c65746f6e}{P36.}] \refstepcounter{paperpremise}\label{premise:50692e696e737453756273696e676c65746f6e}\textbf{Inst subsingleton.} Let \(\alpha\) be a sort. Let \(\beta\) be a function of type \(\alpha \to \mathsf{Sort}\). Let \(\mathit{inst}\) be a proof of \(\forall a\, :\, \alpha,\; \beta\left(a\right)\text{ has at most one element}\). Then \(\left(\Pi_{a\, :\, \alpha}\, \beta\left(a\right)\right)\text{ has at most one element}\).
\item[\paperlineanchor{premise:50692e6f6e655f6170706c79}\paperlink{premise:50692e6f6e655f6170706c79}{P37.}] \refstepcounter{paperpremise}\label{premise:50692e6f6e655f6170706c79}\textbf{One apply.} Let \(I\) be a type. Let \(f\) be a function of type \(I \to \mathsf{Type}\). Let \(i\) be an element of \(I\). Fix a structure \(\mathit{inst}\) of type \(\Pi_{i_{1}\, :\, I}\, \operatorname{One}\left(f\left(i_{1}\right)\right)\). Then \(\operatorname{of\,Nat}\left(1, i\right) = 1\).
\item[\paperlineanchor{premise:5365742e696e64696361746f725f6f665f6d656d}\paperlink{premise:5365742e696e64696361746f725f6f665f6d656d}{P38.}] \refstepcounter{paperpremise}\label{premise:5365742e696e64696361746f725f6f665f6d656d}\textbf{Indicator of membership.} Let \(\alpha\) and \(M\) be types. Fix a structure \(\mathit{inst}\) of type \(\operatorname{Zero}\left(M\right)\). Let \(s\) be a subset of \(\alpha\). Let \(a\) be an element of \(\alpha\). Assume \(a \in s\). Let \(f\) be a function of type \(\alpha \to M\). Then \(\begin{cases}f(a) & a\in s\\ 0 & a\notin s\end{cases} = f\left(a\right)\).
\item[\paperlineanchor{premise:5365742e696e64696361746f725f6f665f6e6f745f6d656d}\paperlink{premise:5365742e696e64696361746f725f6f665f6e6f745f6d656d}{P39.}] \refstepcounter{paperpremise}\label{premise:5365742e696e64696361746f725f6f665f6e6f745f6d656d}\textbf{Indicator of not membership.} Let \(\alpha\) and \(M\) be types. Fix a structure \(\mathit{inst}\) of type \(\operatorname{Zero}\left(M\right)\). Let \(s\) be a subset of \(\alpha\). Let \(a\) be an element of \(\alpha\). Assume \(a \notin s\). Let \(f\) be a function of type \(\alpha \to M\). Then \(\begin{cases}f(a) & a\in s\\ 0 & a\notin s\end{cases} = 0\).
\item[\paperlineanchor{premise:5365742e6d656d5f656d7074795f6966665f66616c7365}\paperlink{premise:5365742e6d656d5f656d7074795f6966665f66616c7365}{P40.}] \refstepcounter{paperpremise}\label{premise:5365742e6d656d5f656d7074795f6966665f66616c7365}\textbf{Membership empty if and only if false.} Let \(\alpha\) be a type. Let \(x\) be an element of \(\alpha\). Then \(x \in \varnothing \iff \mathrm{False}\).
\item[\paperlineanchor{premise:5365742e6d656d5f69496e746572}\paperlink{premise:5365742e6d656d5f69496e746572}{P41.}] \refstepcounter{paperpremise}\label{premise:5365742e6d656d5f69496e746572}\textbf{Membership i inter.} Let \(\alpha\) be a type. Let \(\iota\) be a sort. Let \(x\) be an element of \(\alpha\). Let \(s\) be a function of type \(\iota \to \mathcal{P}(\alpha)\). Then \(x \in \bigcap_{i\, :\, \iota} s\left(i\right) \iff \left(\forall i\, :\, \iota,\; x \in s\left(i\right)\right)\).
\item[\paperlineanchor{premise:53756273696e676c65746f6e2e656c696d}\paperlink{premise:53756273696e676c65746f6e2e656c696d}{P42.}] \refstepcounter{paperpremise}\label{premise:53756273696e676c65746f6e2e656c696d}\textbf{Elim.} Let \(\alpha\) be a sort. Assume \(\alpha\text{ has at most one element}\). Let \(a\) and \(b\) be elements of \(\alpha\). Then \(a = b\).
\item[\paperlineanchor{premise:537562747970652e70726f7065727479}\paperlink{premise:537562747970652e70726f7065727479}{P43.}] \refstepcounter{paperpremise}\label{premise:537562747970652e70726f7065727479}\textbf{Property.} Let \(\alpha\) be a sort. Let \(p\) be a function of type \(\alpha \to \mathsf{Prop}\). Let \(\mathit{self}\) be an element of \(\{x\in \alpha \mid p(x)\}\). Then \(p\left(\mathit{self}\right)\).
\item[\paperlineanchor{premise:6164645f7a65726f}\paperlink{premise:6164645f7a65726f}{P44.}] \refstepcounter{paperpremise}\label{premise:6164645f7a65726f}\textbf{Sum zero.} Let \(M\) be a type. Fix a structure \(\mathit{inst}\) of type \(\operatorname{Add\,Zero\,Class}\left(M\right)\). Let \(a\) be an element of \(M\). Then \(a + 0 = a\).
\item[\paperlineanchor{premise:6469665f6e6567}\paperlink{premise:6469665f6e6567}{P45.}] \refstepcounter{paperpremise}\label{premise:6469665f6e6567}\textbf{Dif neg.} Let \(c\) be a proposition. Let \(h\) be an element of \(\operatorname{Decidable}\left(c\right)\). Let \(\mathit{hnc}\) be a proof of \(\neg c\). Let \(\alpha\) be a sort. Let \(t\) be a function of type \(c \to \alpha\). Let \(e\) be a function of type \(\neg c \to \alpha\). Then \(\begin{cases}t(h_{1}) &\text{given } h_{1}\text{ proving } c\\ e(h_{1}) &\text{given } h_{1}\text{ proving } \neg c\end{cases} = e(\mathit{hnc})\).
\item[\paperlineanchor{premise:6469665f706f73}\paperlink{premise:6469665f706f73}{P46.}] \refstepcounter{paperpremise}\label{premise:6469665f706f73}\textbf{Dif positivity.} Let \(c\) be a proposition. Let \(h\) be an element of \(\operatorname{Decidable}\left(c\right)\). Let \(\mathit{hc}\) be a proof of \(c\). Let \(\alpha\) be a sort. Let \(t\) be a function of type \(c \to \alpha\). Let \(e\) be a function of type \(\neg c \to \alpha\). Then \(\begin{cases}t(h_{1}) &\text{given } h_{1}\text{ proving } c\\ e(h_{1}) &\text{given } h_{1}\text{ proving } \neg c\end{cases} = t(\mathit{hc})\).
\item[\paperlineanchor{premise:65715f66616c7365}\paperlink{premise:65715f66616c7365}{P47.}] \refstepcounter{paperpremise}\label{premise:65715f66616c7365}\textbf{Equals false.} Let \(p\) be a proposition. Assume \(\neg p\). Then \(p = \mathrm{False}\).
\item[\paperlineanchor{premise:65715f6f665f686571}\paperlink{premise:65715f6f665f686571}{P48.}] \refstepcounter{paperpremise}\label{premise:65715f6f665f686571}\textbf{Equals of heq.} Let \(\alpha\) be a sort. Let \(a\) and \(a^{'}\) be elements of \(\alpha\). Assume \(\left(a\, :\, \alpha\right)\text{ is heterogeneously equal to }\left(a^{'}\, :\, \alpha\right)\). Then \(a = a^{'}\).
\item[\paperlineanchor{premise:65715f7375625f6966665f6164645f6571}\paperlink{premise:65715f7375625f6966665f6164645f6571}{P49.}] \refstepcounter{paperpremise}\label{premise:65715f7375625f6966665f6164645f6571}\textbf{Equals difference if and only if sum equals.} Let \(G\) be a type. Fix a structure \(\mathit{inst}\) of type \(\operatorname{Add\,Group}\left(G\right)\). Let \(a\), \(b\), and \(c\) be elements of \(G\). Then \(a = b - c \iff a + c = b\).
\item[\paperlineanchor{premise:666f72616c6c5f636f6e6772}\paperlink{premise:666f72616c6c5f636f6e6772}{P50.}] \refstepcounter{paperpremise}\label{premise:666f72616c6c5f636f6e6772}\textbf{Forall congruence.} Let \(\alpha\) be a sort. Let \(p\) and \(q\) be functions of type \(\alpha \to \mathsf{Prop}\). Assume \(\forall a\, :\, \alpha,\; p\left(a\right) = q\left(a\right)\). Then \((\forall a\, :\, \alpha,\; p(a)) = (\forall a\, :\, \alpha,\; q(a))\).
\item[\paperlineanchor{premise:66756e657874}\paperlink{premise:66756e657874}{P51.}] \refstepcounter{paperpremise}\label{premise:66756e657874}\textbf{Function extensionality.} Let \(\alpha\) be a sort. Let \(\beta\) be a function of type \(\alpha \to \mathsf{Sort}\). Let \(f\) be a function of type \(\Pi_{x\, :\, \alpha}\, \beta\left(x\right)\). Let \(g\) be a function of type \(\Pi_{x\, :\, \alpha}\, \beta\left(x\right)\). Assume \(\forall x\, :\, \alpha,\; f\left(x\right) = g\left(x\right)\). Then \(f = g\).
\item[\paperlineanchor{premise:6966665f616e645f73656c66}\paperlink{premise:6966665f616e645f73656c66}{P52.}] \refstepcounter{paperpremise}\label{premise:6966665f616e645f73656c66}\textbf{If and only if and self.} Let \(p\) and \(q\) be propositions. Then \(\left(p \iff q \land p\right) \iff p \implies q\).
\item[\paperlineanchor{premise:696d706c6965735f636f6e6772}\paperlink{premise:696d706c6965735f636f6e6772}{P53.}] \refstepcounter{paperpremise}\label{premise:696d706c6965735f636f6e6772}\textbf{Implies congruence.} Let \(p_{1}\) and \(p_{2}\) be sorts. Let \(q_{1}\) and \(q_{2}\) be sorts. Assume \(p_{1} = p_{2}\). Assume \(q_{1} = q_{2}\). Then \(\left(p_{1} \to q_{1}\right) = \left(p_{2} \to q_{2}\right)\).
\item[\paperlineanchor{premise:696e737453756273696e676c65746f6e446563696461626c65}\paperlink{premise:696e737453756273696e676c65746f6e446563696461626c65}{P54.}] \refstepcounter{paperpremise}\label{premise:696e737453756273696e676c65746f6e446563696461626c65}\textbf{Inst subsingleton decidable.} Let \(p\) be a proposition. Then \(\operatorname{Decidable}\left(p\right)\text{ has at most one element}\).
\item[\paperlineanchor{premise:6974655f636f6e645f65715f74727565}\paperlink{premise:6974655f636f6e645f65715f74727565}{P55.}] \refstepcounter{paperpremise}\label{premise:6974655f636f6e645f65715f74727565}\textbf{Ite cond equals true.} Let \(\alpha\) be a sort. Let \(c\) be a proposition. Let \(x\) be an element of \(\operatorname{Decidable}\left(c\right)\). Let \(a\) and \(b\) be elements of \(\alpha\). Assume \(c = \mathrm{True}\). Then \(\begin{cases}a &\text{if } c\\ b &\text{otherwise}\end{cases} = a\).
\item[\paperlineanchor{premise:6974655f636f6e6772}\paperlink{premise:6974655f636f6e6772}{P56.}] \refstepcounter{paperpremise}\label{premise:6974655f636f6e6772}\textbf{Ite congruence.} Let \(\alpha\) be a sort. Let \(b\) and \(c\) be propositions. Let \(x\), \(y\), \(u\), and \(v\) be elements of \(\alpha\). Let \(s\) be an element of \(\operatorname{Decidable}\left(b\right)\). Assume \(c\) is decidable. Assume \(b = c\). Assume \(c \implies x = u\). Assume \(\neg c \implies y = v\). Then \(\begin{cases}x &\text{if } b\\ y &\text{otherwise}\end{cases} = \begin{cases}u &\text{if } c\\ v &\text{otherwise}\end{cases}\).
\item[\paperlineanchor{premise:6974655f736d756c}\paperlink{premise:6974655f736d756c}{P57.}] \refstepcounter{paperpremise}\label{premise:6974655f736d756c}\textbf{Ite smul.} Let \(\alpha\) and \(\beta\) be types. Fix a structure \(\mathit{inst}\) of type \(\operatorname{SMul}\left(\beta, \alpha\right)\). Let \(p\) be a proposition. Assume \(p\) is decidable. Let \(a\) be an element of \(\alpha\). Let \(b\) and \(c\) be elements of \(\beta\). Then \(\begin{cases}b &\text{if } p\\ c &\text{otherwise}\end{cases} \mathbin{\bullet} a = \begin{cases}b \mathbin{\bullet} a &\text{if } p\\ c \mathbin{\bullet} a &\text{otherwise}\end{cases}\).
\item[\paperlineanchor{premise:6d756c5f6173736f63}\paperlink{premise:6d756c5f6173736f63}{P58.}] \refstepcounter{paperpremise}\label{premise:6d756c5f6173736f63}\textbf{Multiplication associativity.} Let \(G\) be a type. Fix a structure \(\mathit{inst}\) of type \(\operatorname{Semigroup}\left(G\right)\). Let \(a\), \(b\), and \(c\) be elements of \(G\). Then \(a \cdot b \cdot c = a \cdot \left(b \cdot c\right)\).
\item[\paperlineanchor{premise:6d756c5f697465}\paperlink{premise:6d756c5f697465}{P59.}] \refstepcounter{paperpremise}\label{premise:6d756c5f697465}\textbf{Product ite.} Let \(\alpha\) be a type. Let \(P\) be a proposition. Assume \(P\) is decidable. Fix a structure \(\mathit{inst}_{1}\) of type \(\operatorname{Mul}\left(\alpha\right)\). Let \(a\), \(b\), and \(c\) be elements of \(\alpha\). Then \(a \cdot \begin{cases}b &\text{if } P\\ c &\text{otherwise}\end{cases} = \begin{cases}a \cdot b &\text{if } P\\ a \cdot c &\text{otherwise}\end{cases}\).
\item[\paperlineanchor{premise:6d756c5f6e6567}\paperlink{premise:6d756c5f6e6567}{P60.}] \refstepcounter{paperpremise}\label{premise:6d756c5f6e6567}\textbf{Product neg.} Let \(\alpha\) be a type. Fix a structure \(\mathit{inst}\) of type \(\operatorname{Mul}\left(\alpha\right)\). Fix a structure \(\mathit{inst}_{1}\) of type \(\operatorname{Has\,Distrib\,Neg}\left(\alpha\right)\). Let \(a\) and \(b\) be elements of \(\alpha\). Then \(a \cdot -b = -\left(a \cdot b\right)\).
\item[\paperlineanchor{premise:6d756c5f6f6e65}\paperlink{premise:6d756c5f6f6e65}{P61.}] \refstepcounter{paperpremise}\label{premise:6d756c5f6f6e65}\textbf{Product one.} Let \(M\) be a type. Fix a structure \(\mathit{inst}\) of type \(\operatorname{Mul\,One\,Class}\left(M\right)\). Let \(a\) be an element of \(M\). Then \(a \cdot 1 = a\).
\item[\paperlineanchor{premise:6e65675f6e6567}\paperlink{premise:6e65675f6e6567}{P62.}] \refstepcounter{paperpremise}\label{premise:6e65675f6e6567}\textbf{Neg neg.} Let \(G\) be a type. Fix a structure \(\mathit{inst}\) of type \(\operatorname{Involutive\,Neg}\left(G\right)\). Let \(a\) be an element of \(G\). Then \(-\left(-a\right) = a\).
\item[\paperlineanchor{premise:6e65675f736d756c}\paperlink{premise:6e65675f736d756c}{P63.}] \refstepcounter{paperpremise}\label{premise:6e65675f736d756c}\textbf{Neg smul.} Let \(R\) and \(M\) be types. Fix a structure \(\mathit{inst}\) of type \(\operatorname{Ring}\left(R\right)\). Equip \(M\) with an abelian group structure. Fix a structure \(\mathit{inst}_{2}\) of type \(\operatorname{Module}\left(R, M\right)\). Let \(r\) be an element of \(R\). Let \(x\) be an element of \(M\). Then \(-r \mathbin{\bullet} x = -\left(r \mathbin{\bullet} x\right)\).
\item[\paperlineanchor{premise:6e6f745f66616c73655f65715f74727565}\paperlink{premise:6e6f745f66616c73655f65715f74727565}{P64.}] \refstepcounter{paperpremise}\label{premise:6e6f745f66616c73655f65715f74727565}\textbf{Not false equals true.} \(\left(\neg \mathrm{False}\right) = \mathrm{True}\).
\item[\paperlineanchor{premise:6e736d756c5f65715f6d756c}\paperlink{premise:6e736d756c5f65715f6d756c}{P65.}] \refstepcounter{paperpremise}\label{premise:6e736d756c5f65715f6d756c}\textbf{Repeated addition equals product.} Let \(\alpha\) be a type. Fix a structure \(\mathit{inst}\) of type \(\operatorname{Non\,Assoc\,Semiring}\left(\alpha\right)\). Let \(n\) be a natural number. Let \(a\) be an element of \(\alpha\). Then \(n \mathbin{\bullet} a = \operatorname{cast}_{\alpha}\left(n\right) \cdot a\).
\item[\paperlineanchor{premise:6f6e655f736d756c}\paperlink{premise:6f6e655f736d756c}{P66.}] \refstepcounter{paperpremise}\label{premise:6f6e655f736d756c}\textbf{One smul.} Let \(M\) and \(\alpha\) be types. Fix a structure \(\mathit{inst}\) of type \(\operatorname{Monoid}\left(M\right)\). Fix a structure \(\mathit{inst}_{1}\) of type \(\operatorname{Mul\,Action}\left(M, \alpha\right)\). Let \(b\) be an element of \(\alpha\). Then \(1 \mathbin{\bullet} b = b\).
\item[\paperlineanchor{premise:706f775f73756363}\paperlink{premise:706f775f73756363}{P67.}] \refstepcounter{paperpremise}\label{premise:706f775f73756363}\textbf{Power successor.} Let \(M\) be a type. Fix a structure \(\mathit{inst}\) of type \(\operatorname{Monoid}\left(M\right)\). Let \(a\) be an element of \(M\). Let \(n\) be a natural number. Then \(a^{n + 1} = a^{n} \cdot a\).
\item[\paperlineanchor{premise:706f775f7a65726f}\paperlink{premise:706f775f7a65726f}{P68.}] \refstepcounter{paperpremise}\label{premise:706f775f7a65726f}\textbf{Power zero.} Let \(M\) be a type. Fix a structure \(\mathit{inst}\) of type \(\operatorname{Monoid}\left(M\right)\). Let \(a\) be an element of \(M\). Then \(a^{0} = 1\).
\item[\paperlineanchor{premise:70726f645f696e64696361746f725f6170706c79}\paperlink{premise:70726f645f696e64696361746f725f6170706c79}{P69.}] \refstepcounter{paperpremise}\label{premise:70726f645f696e64696361746f725f6170706c79}\textbf{Product indicator apply.} Let \(\iota\), \(\kappa\), and \(\alpha\) be types. Equip \(\alpha\) with a commutative semiring structure. Let \(s\) be a finite subset of \(\iota\). Let \(f\) be a function of type \(\iota \to \mathcal{P}(\kappa)\). Let \(g\) be a function of type \(\iota \to \kappa \to \alpha\). Let \(j\) be an element of \(\kappa\). Then \(\prod_{i\in s}(\begin{cases}g(i)(j) & j\in f(i)\\ 0 & j\notin f(i)\end{cases}) = \begin{cases}(\prod_{i\in s}(g(i)))(j) & j\in \bigcap_{x\in s} f(x)\\ 0 & j\notin \bigcap_{x\in s} f(x)\end{cases}\).
\item[\paperlineanchor{premise:736d756c5f64697465}\paperlink{premise:736d756c5f64697465}{P70.}] \refstepcounter{paperpremise}\label{premise:736d756c5f64697465}\textbf{Smul dite.} Let \(\alpha\) and \(\beta\) be types. Fix a structure \(\mathit{inst}\) of type \(\operatorname{SMul}\left(\beta, \alpha\right)\). Let \(p\) be a proposition. Assume \(p\) is decidable. Let \(a\) be a function of type \(p \to \alpha\). Let \(b\) be a function of type \(\neg p \to \alpha\). Let \(c\) be an element of \(\beta\). Then \(c \mathbin{\bullet} \begin{cases}a(h) &\text{given } h\text{ proving } p\\ b(h) &\text{given } h\text{ proving } \neg p\end{cases} = \begin{cases}c \mathbin{\bullet} a(h) &\text{given } h\text{ proving } p\\ c \mathbin{\bullet} b(h) &\text{given } h\text{ proving } \neg p\end{cases}\).
\item[\paperlineanchor{premise:7375625f65715f6164645f6e6567}\paperlink{premise:7375625f65715f6164645f6e6567}{P71.}] \refstepcounter{paperpremise}\label{premise:7375625f65715f6164645f6e6567}\textbf{Difference equals sum neg.} Let \(G\) be a type. Fix a structure \(\mathit{inst}\) of type \(\operatorname{Sub\,Neg\,Monoid}\left(G\right)\). Let \(a\) and \(b\) be elements of \(G\). Then \(a - b = a + -b\).
\item[\paperlineanchor{premise:7375625f65715f6e65675f616464}\paperlink{premise:7375625f65715f6e65675f616464}{P72.}] \refstepcounter{paperpremise}\label{premise:7375625f65715f6e65675f616464}\textbf{Difference equals neg sum.} Let \(\alpha\) be a type. Fix a structure \(\mathit{inst}\) of type \(\operatorname{Subtraction\,Comm\,Monoid}\left(\alpha\right)\). Let \(a\) and \(b\) be elements of \(\alpha\). Then \(a - b = -b + a\).
\item[\paperlineanchor{premise:7375625f65715f7a65726f}\paperlink{premise:7375625f65715f7a65726f}{P73.}] \refstepcounter{paperpremise}\label{premise:7375625f65715f7a65726f}\textbf{Difference equals zero.} Let \(G\) be a type. Fix a structure \(\mathit{inst}\) of type \(\operatorname{Add\,Group}\left(G\right)\). Let \(a\) and \(b\) be elements of \(G\). Then \(a - b = 0 \iff a = b\).
\item[\paperlineanchor{premise:7375625f73656c66}\paperlink{premise:7375625f73656c66}{P74.}] \refstepcounter{paperpremise}\label{premise:7375625f73656c66}\textbf{Difference self.} Let \(G\) be a type. Fix a structure \(\mathit{inst}\) of type \(\operatorname{Add\,Group}\left(G\right)\). Let \(a\) be an element of \(G\). Then \(a - a = 0\).
\item[\paperlineanchor{premise:7a65726f5f736d756c}\paperlink{premise:7a65726f5f736d756c}{P75.}] \refstepcounter{paperpremise}\label{premise:7a65726f5f736d756c}\textbf{Zero smul.} Let \(M_{0}\) and \(A\) be types. Fix a structure \(\mathit{inst}\) of type \(\operatorname{Zero}\left(M_{0}\right)\). Fix a structure \(\mathit{inst}_{1}\) of type \(\operatorname{Zero}\left(A\right)\). Fix a structure \(\mathit{inst}_{2}\) of type \(\operatorname{SMul\,With\,Zero}\left(M_{0}, A\right)\). Let \(m\) be an element of \(A\). Then \(0 \mathbin{\bullet} m = 0\).
\item[\paperlineanchor{premise:7a65726f5f737562}\paperlink{premise:7a65726f5f737562}{P76.}] \refstepcounter{paperpremise}\label{premise:7a65726f5f737562}\textbf{Zero difference.} Let \(G\) be a type. Fix a structure \(\mathit{inst}\) of type \(\operatorname{Sub\,Neg\,Monoid}\left(G\right)\). Let \(a\) be an element of \(G\). Then \(0 - a = -a\).
\end{description}

\section{Results}
\begin{lemma}\label{theorem:1}
Let \(\iota\) and \(\alpha\) be types. Assume equality on \(\alpha\) is decidable. Let \(s\) be a finite subset of \(\iota\). Assume \(s\ne\varnothing\). Let \(S\) be a family of finite subsets of \(\alpha\) indexed by \(\iota\). Let \(a\) be an element of \(\alpha\). Then
\begin{center}
\(\displaystyle \prod_{i\in s}(\mathbf{1}_{\bigcup_{i_{1}\in s} S(i_{1})}(a) - \mathbf{1}_{S(i)}(a)) = 0.\)
\end{center}
\end{lemma}

\begin{proof}
Distinguish whether \(a \in \bigcup_{i\in s} S(i)\) holds.
Consider the following case.
\begin{proofstep}
Assume \(a \in \bigcup_{i\in s} S(i)\).
By \paperlink{premise:46696e7365742e6d656d5f6269556e696f6e}{Prerequisite~P\ref{premise:46696e7365742e6d656d5f6269556e696f6e}}, we obtain \(a \in \bigcup_{i\in s} S(i) \iff (\exists a_{1}\, :\, \iota,\; a_{1} \in s \land a \in S(a_{1}))\).
Use the hypothesis \(a \in \bigcup_{i\in s} S(i)\).
By the forward implication, we obtain \(\exists a_{1}\, :\, \iota,\; a_{1} \in s \land a \in S\left(a_{1}\right)\).
Consider the following case.
\begin{proofstep}
Fix \(i\) in \(\iota\).
Assume \(i \in s \land a \in S\left(i\right)\).
Use the hypothesis \(i \in s \land a \in S\left(i\right)\).
Consider the following case.
\begin{proofstep}
Assume \(i \in s\).
Assume \(a \in S\left(i\right)\).
Use the hypothesis \(i \in s\).
Rewriting with \paperlink{premise:5365742e696e64696361746f725f6f665f6d656d}{Prerequisite~P\ref{premise:5365742e696e64696361746f725f6f665f6d656d}}, \paperlink{premise:46696e7365742e6d656d5f636f65}{Prerequisite~P\ref{premise:46696e7365742e6d656d5f636f65}}, \(a \in \bigcup_{i\in s} S(i)\), \(a \in S\left(i\right)\) and \paperlink{premise:7375625f73656c66}{Prerequisite~P\ref{premise:7375625f73656c66}} gives \(\mathbf{1}_{\bigcup_{i_{1}\in s} S(i_{1})}(a) - \mathbf{1}_{S\left(i\right)}(a) = 0\).
By \paperlink{premise:46696e7365742e70726f645f65715f7a65726f}{Prerequisite~P\ref{premise:46696e7365742e70726f645f65715f7a65726f}} with \(M_{0} = \mathbb{Z}\), we obtain \(\prod_{j\in s}\left(\mathbf{1}_{\bigcup_{i_{1}\in s} S(i_{1})}(a) - \mathbf{1}_{S\left(j\right)}(a)\right) = 0\).
\end{proofstep}
Using the two conjuncts, we obtain \(\prod_{i_{1}\in s}(\mathbf{1}_{\bigcup_{i_{2}\in s} S(i_{2})}(a) - \mathbf{1}_{S(i_{1})}(a)) = 0\).
\end{proofstep}
Using the existential hypothesis, we obtain \(\prod_{i\in s}\left(\mathbf{1}_{\bigcup_{i_{1}\in s} S(i_{1})}(a) - \mathbf{1}_{S\left(i\right)}(a)\right) = 0\).
\end{proofstep}
Consider the following case.
\begin{proofstep}
Assume \(a \notin \bigcup_{i\in s} S(i)\).
Use the hypothesis \(s\ne\varnothing\).
Consider the following case.
\begin{proofstep}
Fix \(i\) in \(\iota\).
Assume \(i \in s\).
Assume \(a \in S\left(i\right)\).
Use the hypothesis \(i \in s\).
Use the hypothesis \(a \in S\left(i\right)\).
By \paperlink{premise:46696e7365742e7375627365745f6269556e696f6e5f6f665f6d656d}{Prerequisite~P\ref{premise:46696e7365742e7375627365745f6269556e696f6e5f6f665f6d656d}} with \(u = S\), we obtain \(a \in \bigcup_{i_{1}\in s} S(i_{1})\).
Applying the hypothesis \(a \notin \bigcup_{i\in s} S(i)\), we obtain \(\mathrm{False}\).
Use the hypothesis \(i \in s\).
Rewriting with \paperlink{premise:5365742e696e64696361746f725f6f665f6e6f745f6d656d}{Prerequisite~P\ref{premise:5365742e696e64696361746f725f6f665f6e6f745f6d656d}}, \paperlink{premise:46696e7365742e6d656d5f636f65}{Prerequisite~P\ref{premise:46696e7365742e6d656d5f636f65}}, \paperlink{premise:65715f66616c7365}{Prerequisite~P\ref{premise:65715f66616c7365}} (using \(a \notin \bigcup_{i\in s} S(i)\)), \paperlink{premise:6e6f745f66616c73655f65715f74727565}{Prerequisite~P\ref{premise:6e6f745f66616c73655f65715f74727565}}, \paperlink{premise:65715f66616c7365}{Prerequisite~P\ref{premise:65715f66616c7365}} (using \(a \notin S\left(i\right)\)) and \paperlink{premise:7375625f73656c66}{Prerequisite~P\ref{premise:7375625f73656c66}} gives \(\mathbf{1}_{\bigcup_{i_{1}\in s} S(i_{1})}(a) - \mathbf{1}_{S\left(i\right)}(a) = 0\).
By \paperlink{premise:46696e7365742e70726f645f65715f7a65726f}{Prerequisite~P\ref{premise:46696e7365742e70726f645f65715f7a65726f}} with \(M_{0} = \mathbb{Z}\), we obtain \(\prod_{i_{1}\in s}(\mathbf{1}_{\bigcup_{i_{2}\in s} S(i_{2})}(a) - \mathbf{1}_{S(i_{1})}(a)) = 0\).
\end{proofstep}
Using the existential hypothesis, we obtain \(\prod_{i\in s}\left(\mathbf{1}_{\bigcup_{i_{1}\in s} S(i_{1})}(a) - \mathbf{1}_{S\left(i\right)}(a)\right) = 0\).
\end{proofstep}
The proposition or its negation holds. Hence \(\prod_{i\in s}\left(\mathbf{1}_{\bigcup_{i_{1}\in s} S(i_{1})}(a) - \mathbf{1}_{S\left(i\right)}(a)\right) = 0\).
\end{proof}

\begin{lemma}[{Inclusion-exclusion principle for the sum of a function over a union}]\label{theorem:2}
Let \(\iota\), \(\alpha\), and \(G\) be types. Assume equality on \(\alpha\) is decidable. Equip \(G\) with an abelian group structure. Let \(s\) be a finite subset of \(\iota\). Let \(S\) be a family of finite subsets of \(\alpha\) indexed by \(\iota\). Let \(f\) be a function of type \(\alpha \to G\). Then
\begin{center}
\(\displaystyle \sum_{a\in \bigcup_{i\in s} S(i)}(f(a)) = \sum_{\varnothing \ne t \subseteq s}((-1)^{|t| + 1} \mathbin{\bullet} \sum_{a\in \bigcap_{i\in t} S(i)}(f(a))).\)
\end{center}
\end{lemma}

\begin{proof}
\par\smallskip\noindent\emph{Equality 1.}\enspace
The substitution uses the following argument.
\begin{proofstep}
Fix \(x\) in \(\{x_{1}\in \mathcal P(s) \mid x_{1}\ne\varnothing\}\).
Assume \(x \in \operatorname{attach}\left(\{x_{1}\in \mathcal P(s) \mid x_{1}\ne\varnothing\}\right)\).
By \paperlink{premise:706f775f73756363}{Prerequisite~P\ref{premise:706f775f73756363}} with \(M = \mathbb{Z}\) and \(-1\) in place of its parameter \(a\) and \(n = \left|\operatorname{val}(x)\right|\), we obtain \((-1)^{|\operatorname{val}(x)| + 1} = (-1)^{|\operatorname{val}(x)|} \cdot -1\).
By \paperlink{premise:6d756c5f6e6567}{Prerequisite~P\ref{premise:6d756c5f6e6567}} with \(\mathbb{Z}\) as its type parameter \(\alpha\) and \(\left(-1\right)^{\left|\operatorname{val}(x)\right|}\) in place of its parameter \(a\) and \(b = 1\), we obtain \((-1)^{|\operatorname{val}(x)|} \cdot -1 = -((-1)^{|\operatorname{val}(x)|} \cdot 1)\).
By \paperlink{premise:6d756c5f6f6e65}{Prerequisite~P\ref{premise:6d756c5f6f6e65}} with \(M = \mathbb{Z}\) and \(\left(-1\right)^{\left|\operatorname{val}(x)\right|}\) in place of its parameter \(a\), we obtain \((-1)^{|\operatorname{val}(x)|} \cdot 1 = (-1)^{|\operatorname{val}(x)|}\).
Applying the same function to equal arguments, we obtain \(-((-1)^{|\operatorname{val}(x)|} \cdot 1) = -(-1)^{|\operatorname{val}(x)|}\).
Chaining the equalities, we obtain \((-1)^{|\operatorname{val}(x)|} \cdot -1 = -(-1)^{|\operatorname{val}(x)|}\).
Chaining the equalities, we obtain \((-1)^{|\operatorname{val}(x)| + 1} = -(-1)^{|\operatorname{val}(x)|}\).
Applying the same function to equal arguments, we obtain \((-1)^{|\operatorname{val}(x)| + 1} \mathbin{\bullet} \sum_{a_{1}\in \bigcap_{i\in \operatorname{val}(x)} S(i)}(f(a_{1})) = -(-1)^{|\operatorname{val}(x)|} \mathbin{\bullet} \sum_{a_{1}\in \bigcap_{i\in \operatorname{val}(x)} S(i)}(f(a_{1}))\).
By \paperlink{premise:6e65675f736d756c}{Prerequisite~P\ref{premise:6e65675f736d756c}} with \(R = \mathbb{Z}\) and \(r = \left(-1\right)^{\left|\operatorname{val}(x)\right|}\) and \(\sum_{a_{1}\in \bigcap_{i\in \operatorname{val}(x)} S(i)}\left(f\left(a_{1}\right)\right)\) in place of its parameter \(x\), we obtain \(-(-1)^{|\operatorname{val}(x)|} \mathbin{\bullet} \sum_{a_{1}\in \bigcap_{i\in \operatorname{val}(x)} S(i)}(f(a_{1})) = -((-1)^{|\operatorname{val}(x)|} \mathbin{\bullet} \sum_{a_{1}\in \bigcap_{i\in \operatorname{val}(x)} S(i)}(f(a_{1})))\).
Chaining the equalities, we obtain \((-1)^{|\operatorname{val}(x)| + 1} \mathbin{\bullet} \sum_{a_{1}\in \bigcap_{i\in \operatorname{val}(x)} S(i)}(f(a_{1})) = -((-1)^{|\operatorname{val}(x)|} \mathbin{\bullet} \sum_{a_{1}\in \bigcap_{i\in \operatorname{val}(x)} S(i)}(f(a_{1})))\).
\end{proofstep}
The substitution uses the following argument.
\begin{proofstep}
Fix \(x\) in \(\{x_{1}\in \mathcal P(s) \mid x_{1}\ne\varnothing\}\).
Assume \(x \in \operatorname{attach}\left(\{x_{1}\in \mathcal P(s) \mid x_{1}\ne\varnothing\}\right)\).
By \paperlink{premise:6e65675f6e6567}{Prerequisite~P\ref{premise:6e65675f6e6567}} with \((-1)^{|\operatorname{val}(x)|} \mathbin{\bullet} \sum_{a_{1}\in \bigcap_{i\in \operatorname{val}(x)} S(i)}(f(a_{1}))\) in place of its parameter \(a\), we obtain \(-(-((-1)^{|\operatorname{val}(x)|} \mathbin{\bullet} \sum_{a_{1}\in \bigcap_{i\in \operatorname{val}(x)} S(i)}(f(a_{1})))) = (-1)^{|\operatorname{val}(x)|} \mathbin{\bullet} \sum_{a_{1}\in \bigcap_{i\in \operatorname{val}(x)} S(i)}(f(a_{1}))\).
\end{proofstep}
The substitution uses the following argument.
\begin{proofstep}
Fix \(\mathit{inst}_{2}\) in \(\operatorname{Decidable\,Pred}\left(x \mapsto \neg \left(x\ne\varnothing\right)\right)\).
For the required implication, argue as follows.
\begin{proofstep}
Fix \(s_{1}\) in \(\mathcal{P}_{\mathrm{fin}}(\mathcal{P}_{\mathrm{fin}}(\iota))\).
Fix \(s_{2}\) in \(\mathcal{P}_{\mathrm{fin}}(\mathcal{P}_{\mathrm{fin}}(\iota))\).
Assume \(s_{1} = s_{2}\).
By definition, we obtain \(\{x\in s_{1} \mid \neg \left(x\ne\varnothing\right)\} = \{x\in s_{1} \mid \neg \left(x\ne\varnothing\right)\}\).
Use the hypothesis \(s_{1} = s_{2}\).
Substituting equals, we obtain \(\{x\in s_{1} \mid \neg \left(x\ne\varnothing\right)\} = \{x\in s_{2} \mid \neg \left(x\ne\varnothing\right)\}\).
\end{proofstep}
For the required implication, argue as follows.
\begin{proofstep}
Fix \(a\) in \(\mathcal{P}_{\mathrm{fin}}(\iota)\).
By \paperlink{premise:696e737453756273696e676c65746f6e446563696461626c65}{Prerequisite~P\ref{premise:696e737453756273696e676c65746f6e446563696461626c65}} with \(p = \left(x \mapsto \neg \left(x\ne\varnothing\right)\right)\left(a\right)\), we obtain \(\operatorname{Decidable}((x \mapsto \neg (x\ne\varnothing))(a))\text{ has at most one element}\).
\end{proofstep}
By \paperlink{premise:50692e696e737453756273696e676c65746f6e}{Prerequisite~P\ref{premise:50692e696e737453756273696e676c65746f6e}} with \(\mathcal{P}_{\mathrm{fin}}(\iota)\) as its type parameter \(\alpha\), we obtain \((\Pi_{a\, :\, \mathcal{P}_{\mathrm{fin}}(\iota)}\, (a_{1} \mapsto \operatorname{Decidable}((x \mapsto \neg (x\ne\varnothing))(a_{1})))(a))\text{ has at most one element}\).
By \paperlink{premise:53756273696e676c65746f6e2e656c696d}{Prerequisite~P\ref{premise:53756273696e676c65746f6e2e656c696d}} with \(\operatorname{Decidable\,Pred}\left(x \mapsto \neg \left(x\ne\varnothing\right)\right)\) as its type parameter \(\alpha\) and \(a = \mathit{inst}_{2}\) and \(b = \mathit{inst}_{3}\), we obtain \(\mathit{inst}_{2} = \mathit{inst}_{3}\).
Substituting equals, we obtain \((\mathit{inst}_{4} \mapsto \forall s_{1}\, :\, \mathcal{P}_{\mathrm{fin}}(\mathcal{P}_{\mathrm{fin}}(\iota)),\; \forall s_{2}\, :\, \mathcal{P}_{\mathrm{fin}}(\mathcal{P}_{\mathrm{fin}}(\iota)),\; s_{1} = s_{2} \implies \{x\in s_{1} \mid \neg (x\ne\varnothing)\} = \{x\in s_{2} \mid \neg (x\ne\varnothing)\})(\mathit{inst}_{3})\).
\end{proofstep}
The substitution uses the following argument.
\begin{proofstep}
Fix \(x\) in \(\mathcal{P}_{\mathrm{fin}}(\iota)\).
By \paperlink{premise:46696e7365742e6e6f745f6e6f6e656d7074795f6966665f65715f656d707479}{Prerequisite~P\ref{premise:46696e7365742e6e6f745f6e6f6e656d7074795f6966665f65715f656d707479}}, we obtain \(\neg \left(x\ne\varnothing\right) \iff x = \varnothing\).
By propositional extensionality, we obtain \(\left(\neg \left(x\ne\varnothing\right)\right) = \left(x = \varnothing\right)\).
\end{proofstep}
The substitution uses the following argument.
\begin{proofstep}
Fix \(x\) in \(\mathcal{P}_{\mathrm{fin}}(\iota)\).
Assume \(x \in \{\varnothing\}\).
By definition, we obtain \((-1)^{|x|} \mathbin{\bullet} \sum_{a_{1}\in \bigcup_{i\in s} S(i)}(f(a_{1})) = (-1)^{|x|} \mathbin{\bullet} \sum_{a_{1}\in \bigcup_{i\in s} S(i)}(f(a_{1}))\).
\end{proofstep}
Rewriting with \paperlink{premise:46696e7365742e73756d5f636f6e6772}{Prerequisite~P\ref{premise:46696e7365742e73756d5f636f6e6772}}, \paperlink{premise:7375625f65715f6e65675f616464}{Prerequisite~P\ref{premise:7375625f65715f6e65675f616464}}, \paperlink{premise:46696e7365742e73756d5f6e65675f64697374726962}{Prerequisite~P\ref{premise:46696e7365742e73756d5f6e65675f64697374726962}}, \paperlink{premise:66756e657874}{\emph{Function extensionality}} (\paperlink{premise:66756e657874}{Prerequisite~P\ref{premise:66756e657874}}), \paperlink{premise:46696e7365742e66696c7465725f657127}{Prerequisite~P\ref{premise:46696e7365742e66696c7465725f657127}}, \paperlink{premise:6974655f636f6e645f65715f74727565}{Prerequisite~P\ref{premise:6974655f636f6e645f65715f74727565}}, \paperlink{premise:46696e7365742e6d656d5f706f776572736574}{Prerequisite~P\ref{premise:46696e7365742e6d656d5f706f776572736574}}, \paperlink{premise:46696e7365742e656d7074795f737562736574}{Prerequisite~P\ref{premise:46696e7365742e656d7074795f737562736574}}, \paperlink{premise:46696e7365742e73756d5f73696e676c65746f6e}{Prerequisite~P\ref{premise:46696e7365742e73756d5f73696e676c65746f6e}}, \paperlink{premise:706f775f7a65726f}{Prerequisite~P\ref{premise:706f775f7a65726f}}, \paperlink{premise:6f6e655f736d756c}{Prerequisite~P\ref{premise:6f6e655f736d756c}}, the 4 identities established above and the auxiliary result established above gives \(\sum_{a\in \bigcup_{i\in s} S(i)}(f(a)) - \sum_{\varnothing \ne t \subseteq s}((-1)^{|t| + 1} \mathbin{\bullet} \sum_{a\in \bigcap_{i\in t} S(i)}(f(a))) = \sum_{t\in \operatorname{attach}(\{x\in \mathcal P(s) \mid x\ne\varnothing\})}((-1)^{|t|} \mathbin{\bullet} \sum_{a\in \bigcap_{i\in t} S(i)}(f(a))) + \sum_{t\in \{x\in \mathcal P(s) \mid \neg (x\ne\varnothing)\}}((-1)^{|t|} \mathbin{\bullet} \sum_{a\in \bigcup_{i\in s} S(i)}(f(a)))\).
\par\smallskip\noindent\emph{Equality 2.}\enspace
The substitution uses the following argument.
\begin{proofstep}
Fix \(x\) in \(\mathcal{P}_{\mathrm{fin}}(\iota)\).
Assume \(x \in \mathcal P(s)\).
By \paperlink{premise:736d756c5f64697465}{Prerequisite~P\ref{premise:736d756c5f64697465}} with \(\beta = \mathbb{Z}\) and \(p = x\ne\varnothing\) and \(h \mapsto \sum_{a_{1}\in \bigcap_{i\in x} S(i)}\left(f\left(a_{1}\right)\right)\) in place of its parameter \(a\) and \(b = h \mapsto \sum_{a_{1}\in \bigcup_{i\in s} S(i)}\left(f\left(a_{1}\right)\right)\) and \(c = \left(-1\right)^{\left|x\right|}\), we obtain \((-1)^{|x|} \mathbin{\bullet} \begin{cases}\sum_{a_{1}\in \bigcap_{i\in x} S(i)}(f(a_{1})) &\text{given } h\text{ proving } x\ne\varnothing\\ \sum_{a_{1}\in \bigcup_{i\in s} S(i)}(f(a_{1})) &\text{if } \neg (x\ne\varnothing)\end{cases} = \begin{cases}(-1)^{|x|} \mathbin{\bullet} \sum_{a_{1}\in \bigcap_{i\in x} S(i)}(f(a_{1})) &\text{given } h\text{ proving } x\ne\varnothing\\ (-1)^{|x|} \mathbin{\bullet} \sum_{a_{1}\in \bigcup_{i\in s} S(i)}(f(a_{1})) &\text{if } \neg (x\ne\varnothing)\end{cases}\).
\end{proofstep}
Rewriting with \paperlink{premise:46696e7365742e73756d5f636f6e6772}{Prerequisite~P\ref{premise:46696e7365742e73756d5f636f6e6772}}, \paperlink{premise:46696e7365742e73756d5f64697465}{Prerequisite~P\ref{premise:46696e7365742e73756d5f64697465}} and the identity established above gives \(\sum_{\varnothing \ne t \subseteq s}((-1)^{|t|} \mathbin{\bullet} \sum_{a\in \bigcap_{i\in t} S(i)}(f(a))) + \sum_{x\in \operatorname{attach}(\{x\in \mathcal P(s) \mid \neg (x\ne\varnothing)\})}((-1)^{|x|} \mathbin{\bullet} \sum_{a\in \bigcup_{i\in s} S(i)}(f(a))) = \sum_{t\in \mathcal P(s)}((-1)^{|t|} \mathbin{\bullet} \begin{cases}\sum_{a\in \bigcap_{i\in t} S(i)}(f(a)) &\text{given } \mathit{ht}\text{ proving } t\ne\varnothing\\ \sum_{a\in \bigcup_{i\in s} S(i)}(f(a)) &\text{if } \neg (t\ne\varnothing)\end{cases})\).
Rewriting with \paperlink{premise:46696e7365742e73756d5f617474616368}{Prerequisite~P\ref{premise:46696e7365742e73756d5f617474616368}} in the preceding result gives \(\sum_{\varnothing \ne t \subseteq s}((-1)^{|t|} \mathbin{\bullet} \sum_{a\in \bigcap_{i\in t} S(i)}(f(a))) + \sum_{t\in \{x\in \mathcal P(s) \mid \neg (x\ne\varnothing)\}}((-1)^{|t|} \mathbin{\bullet} \sum_{a\in \bigcup_{i\in s} S(i)}(f(a))) = \sum_{t\in \mathcal P(s)}((-1)^{|t|} \mathbin{\bullet} \begin{cases}\sum_{a\in \bigcap_{i\in t} S(i)}(f(a)) &\text{given } \mathit{ht}\text{ proving } t\ne\varnothing\\ \sum_{a\in \bigcup_{i\in s} S(i)}(f(a)) &\text{if } \neg (t\ne\varnothing)\end{cases})\).
\par\smallskip\noindent\emph{Equality 3.}\enspace
The substitution uses the following argument.
\begin{proofstep}
Fix \(x\) in \(\alpha\).
Assume \(x \in \bigcup_{i\in s} S(i)\).
By \paperlink{premise:46696e7365742e70726f645f737562}{Prerequisite~P\ref{premise:46696e7365742e70726f645f737562}} with \(\mathbb{Z}\) as its type parameter \(\alpha\) and \(i \mapsto 1\) in place of its parameter \(f\) and \(g = i \mapsto \mathbf{1}_{S\left(i\right)}(x)\), we obtain \(\prod_{i\in s}(1 - \mathbf{1}_{S(i)}(x)) = \sum_{t\in \mathcal P(s)}((-1)^{|t|} \cdot \prod_{i\in s \setminus t}(1) \cdot \prod_{i\in t}(\mathbf{1}_{S(i)}(x)))\).
By definition, we obtain \(\mathcal P(s) = \mathcal P(s)\).
For the required implication, argue as follows.
\begin{proofstep}
Fix \(x_{1}\) in \(\mathcal{P}_{\mathrm{fin}}(\iota)\).
Assume \(x_{1} \in \mathcal P(s)\).
By \paperlink{premise:6d756c5f6173736f63}{Prerequisite~P\ref{premise:6d756c5f6173736f63}} with \(\mathbb{Z}\) as its type parameter \(G\) and \(\left(-1\right)^{\left|x_{1}\right|}\) in place of its parameter \(a\) and \(b = \prod_{i\in s \setminus x_{1}}\left(1\right)\) and \(c = \prod_{i\in x_{1}}\left(\mathbf{1}_{S\left(i\right)}(x)\right)\), we obtain \((-1)^{|x_{1}|} \cdot \prod_{i\in s \setminus x_{1}}(1) \cdot \prod_{i\in x_{1}}(\mathbf{1}_{S(i)}(x)) = (-1)^{|x_{1}|} \cdot (\prod_{i\in s \setminus x_{1}}(1) \cdot \prod_{i\in x_{1}}(\mathbf{1}_{S(i)}(x)))\).
\end{proofstep}
By \paperlink{premise:46696e7365742e73756d5f636f6e6772}{Prerequisite~P\ref{premise:46696e7365742e73756d5f636f6e6772}} with \(\mathcal{P}_{\mathrm{fin}}(\iota)\) as its type parameter \(\iota\) and \(M = \mathbb{Z}\), we obtain \(\sum_{t\in \mathcal P(s)}((-1)^{|t|} \cdot \prod_{i\in s \setminus t}(1) \cdot \prod_{i\in t}(\mathbf{1}_{S(i)}(x))) = \sum_{x_{1}\in \mathcal P(s)}((-1)^{|x_{1}|} \cdot (\prod_{i\in s \setminus x_{1}}(1) \cdot \prod_{i\in x_{1}}(\mathbf{1}_{S(i)}(x))))\).
Chaining the equalities, we obtain \(\prod_{i\in s}(1 - \mathbf{1}_{S(i)}(x)) = \sum_{x_{1}\in \mathcal P(s)}((-1)^{|x_{1}|} \cdot (\prod_{i\in s \setminus x_{1}}(1) \cdot \prod_{i\in x_{1}}(\mathbf{1}_{S(i)}(x))))\).
Applying the same function to equal arguments, we obtain \(\prod_{i\in s}(1 - \mathbf{1}_{S(i)}(x)) \mathbin{\bullet} f(x) = \sum_{x_{1}\in \mathcal P(s)}((-1)^{|x_{1}|} \cdot (\prod_{i\in s \setminus x_{1}}(1) \cdot \prod_{i\in x_{1}}(\mathbf{1}_{S(i)}(x)))) \mathbin{\bullet} f(x)\).
By \paperlink{premise:46696e7365742e73756d5f736d756c}{Prerequisite~P\ref{premise:46696e7365742e73756d5f736d756c}} with \(\mathcal{P}_{\mathrm{fin}}(\iota)\) as its type parameter \(\iota\) and \(R = \mathbb{Z}\), we obtain \(\sum_{i\in \mathcal P(s)}((-1)^{|i|} \cdot (\prod_{i_{1}\in s \setminus i}(1) \cdot \prod_{i_{1}\in i}(\mathbf{1}_{S(i_{1})}(x)))) \mathbin{\bullet} f(x) = \sum_{i\in \mathcal P(s)}((-1)^{|i|} \cdot (\prod_{i_{1}\in s \setminus i}(1) \cdot \prod_{i_{1}\in i}(\mathbf{1}_{S(i_{1})}(x))) \mathbin{\bullet} f(x))\).
Chaining the equalities, we obtain \(\prod_{i\in s}(1 - \mathbf{1}_{S(i)}(x)) \mathbin{\bullet} f(x) = \sum_{i\in \mathcal P(s)}((-1)^{|i|} \cdot (\prod_{i_{1}\in s \setminus i}(1) \cdot \prod_{i_{1}\in i}(\mathbf{1}_{S(i_{1})}(x))) \mathbin{\bullet} f(x))\).
\end{proofstep}
By definition, we obtain \(\mathcal P(s) = \mathcal P(s)\).
For the required implication, argue as follows.
\begin{proofstep}
Fix \(t\) in \(\mathcal{P}_{\mathrm{fin}}(\iota)\).
Assume \(t \in \mathcal P(s)\).
Distinguish whether \(t\ne\varnothing\) holds.
Consider the following case.
\begin{proofstep}
Assume \(t\ne\varnothing\).
Use the hypothesis \(t\ne\varnothing\).
Consider the following case.
\begin{proofstep}
Fix \(i\) in \(\iota\).
Assume \(i \in t\).
The substitution uses the following argument.
\begin{proofstep}
Fix \(x\) in \(\alpha\).
Assume \(x \in \bigcup_{i_{1}\in s} S(i_{1})\).
By \paperlink{premise:46696e7365742e70726f645f636f6e73745f6f6e65}{Prerequisite~P\ref{premise:46696e7365742e70726f645f636f6e73745f6f6e65}} with \(\beta = \mathbb{Z}\), we obtain \(\prod_{x_{1}\in s \setminus t}\left(1\right) = 1\).
Applying the same function to equal arguments, we obtain \(\left(y \mapsto \prod_{x_{1}\in s \setminus t}\left(1\right) \cdot y\right) = \left(y \mapsto 1 \cdot y\right)\).
By \paperlink{premise:70726f645f696e64696361746f725f6170706c79}{Prerequisite~P\ref{premise:70726f645f696e64696361746f725f6170706c79}} with \(\mathbb{Z}\) as its type parameter \(\alpha\) and \(t\) in place of its parameter \(s\) and \(i_{1} \mapsto \operatorname{underlying\,set}(S\left(i_{1}\right))\) in place of its parameter \(f\) and \(g = i_{1} \mapsto 1\) and \(j = x\), we obtain \(\prod_{i_{1}\in t}(\mathbf{1}_{S(i_{1})}(x)) = \begin{cases}(\prod_{i_{1}\in t}(1))(x) & x\in \bigcap_{x_{1}\in t} S(x_{1})\\ 0 & x\notin \bigcap_{x_{1}\in t} S(x_{1})\end{cases}\).
By \paperlink{premise:46696e7365742e70726f645f636f6e73745f6f6e65}{Prerequisite~P\ref{premise:46696e7365742e70726f645f636f6e73745f6f6e65}} with \(\beta = \alpha \to \mathbb{Z}\), we obtain \(\prod_{x_{1}\in t}\left(1\right) = 1\).
Applying the same function to equal arguments, we obtain \(\begin{cases}(\prod_{x_{1}\in t}(1))(x) & x\in \bigcap_{x_{1}\in t} S(x_{1})\\ 0 & x\notin \bigcap_{x_{1}\in t} S(x_{1})\end{cases} = \mathbf{1}_{\bigcap_{x_{1}\in t} S(x_{1})}(x)\).
Chaining the equalities, we obtain \(\prod_{i_{1}\in t}(\mathbf{1}_{S(i_{1})}(x)) = \mathbf{1}_{\bigcap_{x_{1}\in t} S(x_{1})}(x)\).
By congruence, we obtain \(\prod_{x_{1}\in s \setminus t}(1) \cdot \prod_{i_{1}\in t}(\mathbf{1}_{S(i_{1})}(x)) = 1 \cdot \mathbf{1}_{\bigcap_{x_{1}\in t} S(x_{1})}(x)\).
Applying the same function to equal arguments, we obtain \((-1)^{|t|} \cdot (\prod_{x_{1}\in s \setminus t}(1) \cdot \prod_{i_{1}\in t}(\mathbf{1}_{S(i_{1})}(x))) \mathbin{\bullet} f(x) = (-1)^{|t|} \cdot (1 \cdot \mathbf{1}_{\bigcap_{x_{1}\in t} S(x_{1})}(x)) \mathbin{\bullet} f(x)\).
\end{proofstep}
The substitution uses the following argument.
\begin{proofstep}
Fix \(x\) in \(\alpha\).
Assume \(x \in \bigcup_{i_{1}\in s} S(i_{1})\).
By \paperlink{premise:5365742e6d656d5f69496e746572}{Prerequisite~P\ref{premise:5365742e6d656d5f69496e746572}}, we obtain \(x \in \bigcap_{i_{1}\, :\, \iota} (i_{2} \mapsto \bigcap_{h\text{ proving }i_{2} \in t} S(i_{2}))(i_{1}) \iff (\forall i_{1}\, :\, \iota,\; x \in (i_{2} \mapsto \bigcap_{h\text{ proving }i_{2} \in t} S(i_{2}))(i_{1}))\).
By propositional extensionality, we obtain \((x \in \bigcap_{i_{1}\in t} S(i_{1})) = (\forall i_{1}\, :\, \iota,\; x \in \bigcap_{h\text{ proving }i_{1} \in t} S(i_{1}))\).
For the required implication, argue as follows.
\begin{proofstep}
Fix \(i_{1}\) in \(\iota\).
By \paperlink{premise:5365742e6d656d5f69496e746572}{Prerequisite~P\ref{premise:5365742e6d656d5f69496e746572}} with \(i_{1} \in t\) as its type parameter \(\iota\), we obtain \(x \in \bigcap_{i_{2}\, :\, i_{1} \in t} (i_{3} \mapsto \operatorname{underlying\,set}(S(i_{1})))(i_{2}) \iff (\forall i_{2}\, :\, i_{1} \in t,\; x \in (i_{3} \mapsto \operatorname{underlying\,set}(S(i_{1})))(i_{2}))\).
By propositional extensionality, we obtain \((x \in \bigcap_{i_{2}\text{ proving }i_{1} \in t} S(i_{1})) = (i_{1} \in t \implies x \in \operatorname{underlying\,set}(S(i_{1})))\).
By definition, we obtain \(\left(i_{1} \in t\right) = \left(i_{1} \in t\right)\).
By \paperlink{premise:46696e7365742e6d656d5f636f65}{Prerequisite~P\ref{premise:46696e7365742e6d656d5f636f65}}, we obtain \(x \in \operatorname{underlying\,set}(S\left(i_{1}\right)) \iff x \in S\left(i_{1}\right)\).
By propositional extensionality, we obtain \((x \in \operatorname{underlying\,set}(S(i_{1}))) = (x \in S(i_{1}))\).
By \paperlink{premise:696d706c6965735f636f6e6772}{Prerequisite~P\ref{premise:696d706c6965735f636f6e6772}} with \(p_{1} = i_{1} \in t\) and \(p_{2} = i_{1} \in t\) and \(q_{1} = x \in \operatorname{underlying\,set}(S\left(i_{1}\right))\) and \(q_{2} = x \in S\left(i_{1}\right)\), we obtain \((i_{1} \in t \implies x \in \operatorname{underlying\,set}(S(i_{1}))) = (i_{1} \in t \implies x \in S(i_{1}))\).
Chaining the equalities, we obtain \((x \in \bigcap_{i_{2}\text{ proving }i_{1} \in t} S(i_{1})) = (i_{1} \in t \implies x \in S(i_{1}))\).
\end{proofstep}
By \paperlink{premise:666f72616c6c5f636f6e6772}{Prerequisite~P\ref{premise:666f72616c6c5f636f6e6772}}, we obtain \((\forall a_{2}\, :\, \iota,\; x \in \bigcap_{i_{1}\text{ proving }a_{2} \in t} S(a_{2})) = (\forall a_{2}\, :\, \iota,\; a_{2} \in t \implies x \in S(a_{2}))\).
Chaining the equalities, we obtain \((x \in \bigcap_{i_{1}\in t} S(i_{1})) = (\forall a_{2}\, :\, \iota,\; a_{2} \in t \implies x \in S(a_{2}))\).
For the required implication, argue as follows.
\begin{proofstep}
Assume \(\forall i_{1}\, :\, \iota,\; i_{1} \in t \implies x \in S\left(i_{1}\right)\).
By definition, we obtain \(1 = 1\).
\end{proofstep}
For the required implication, argue as follows.
\begin{proofstep}
Assume \(\neg \left(\forall i_{1}\, :\, \iota,\; i_{1} \in t \implies x \in S\left(i_{1}\right)\right)\).
By definition, we obtain \(0 = 0\).
\end{proofstep}
By \paperlink{premise:6974655f636f6e6772}{Prerequisite~P\ref{premise:6974655f636f6e6772}} with \(\mathbb{Z}\) as its type parameter \(\alpha\), we obtain \(\begin{cases}1 &\text{if } x \in \bigcap_{x_{1}\in t} S(x_{1})\\ 0 &\text{otherwise}\end{cases} = \begin{cases}1 &\text{if } \forall i_{1}\, :\, \iota,\; i_{1} \in t \implies x \in S(i_{1})\\ 0 &\text{otherwise}\end{cases}\).
Applying the same function to equal arguments, we obtain \(1 \cdot \begin{cases}1 &\text{if } x \in \bigcap_{x_{1}\in t} S(x_{1})\\ 0 &\text{otherwise}\end{cases} = 1 \cdot \begin{cases}1 &\text{if } \forall i_{1}\, :\, \iota,\; i_{1} \in t \implies x \in S(i_{1})\\ 0 &\text{otherwise}\end{cases}\).
By \paperlink{premise:6d756c5f697465}{Prerequisite~P\ref{premise:6d756c5f697465}} with \(\mathbb{Z}\) as its type parameter \(\alpha\) and \(P = \forall i_{1}\, :\, \iota,\; i_{1} \in t \implies x \in S\left(i_{1}\right)\) and \(1\) in place of its parameter \(a\) and \(b = 1\) and \(c = 0\), we obtain \(1 \cdot \begin{cases}1 &\text{if } \forall i_{1}\, :\, \iota,\; i_{1} \in t \implies x \in S(i_{1})\\ 0 &\text{otherwise}\end{cases} = \begin{cases}1 \cdot 1 &\text{if } \forall i_{1}\, :\, \iota,\; i_{1} \in t \implies x \in S(i_{1})\\ 1 \cdot 0 &\text{otherwise}\end{cases}\).
Chaining the equalities, we obtain \(1 \cdot \begin{cases}1 &\text{if } x \in \bigcap_{x_{1}\in t} S(x_{1})\\ 0 &\text{otherwise}\end{cases} = \begin{cases}1 \cdot 1 &\text{if } \forall i_{1}\, :\, \iota,\; i_{1} \in t \implies x \in S(i_{1})\\ 1 \cdot 0 &\text{otherwise}\end{cases}\).
By definition, we obtain \((\forall i_{1}\, :\, \iota,\; i_{1} \in t \implies x \in S(i_{1})) = (\forall i_{1}\, :\, \iota,\; i_{1} \in t \implies x \in S(i_{1}))\).
For the required implication, argue as follows.
\begin{proofstep}
Assume \(\forall i_{1}\, :\, \iota,\; i_{1} \in t \implies x \in S\left(i_{1}\right)\).
By \paperlink{premise:6d756c5f6f6e65}{Prerequisite~P\ref{premise:6d756c5f6f6e65}} with \(M = \mathbb{Z}\) and \(1\) in place of its parameter \(a\), we obtain \(1 \cdot 1 = 1\).
\end{proofstep}
For the required implication, argue as follows.
\begin{proofstep}
Assume \(\neg \left(\forall i_{1}\, :\, \iota,\; i_{1} \in t \implies x \in S\left(i_{1}\right)\right)\).
By \paperlink{premise:4d756c5a65726f436c6173732e6d756c5f7a65726f}{Prerequisite~P\ref{premise:4d756c5a65726f436c6173732e6d756c5f7a65726f}} with \(M_{0} = \mathbb{Z}\) and \(1\) in place of its parameter \(a\), we obtain \(1 \cdot 0 = 0\).
\end{proofstep}
By \paperlink{premise:6974655f636f6e6772}{Prerequisite~P\ref{premise:6974655f636f6e6772}} with \(\mathbb{Z}\) as its type parameter \(\alpha\), we obtain \(\begin{cases}1 \cdot 1 &\text{if } \forall i_{1}\, :\, \iota,\; i_{1} \in t \implies x \in S(i_{1})\\ 1 \cdot 0 &\text{otherwise}\end{cases} = \begin{cases}1 &\text{if } \forall i_{1}\, :\, \iota,\; i_{1} \in t \implies x \in S(i_{1})\\ 0 &\text{otherwise}\end{cases}\).
Chaining the equalities, we obtain \(1 \cdot \begin{cases}1 &\text{if } x \in \bigcap_{x_{1}\in t} S(x_{1})\\ 0 &\text{otherwise}\end{cases} = \begin{cases}1 &\text{if } \forall i_{1}\, :\, \iota,\; i_{1} \in t \implies x \in S(i_{1})\\ 0 &\text{otherwise}\end{cases}\).
Applying the same function to equal arguments, we obtain \((-1)^{|t|} \cdot (1 \cdot \begin{cases}1 &\text{if } x \in \bigcap_{x_{1}\in t} S(x_{1})\\ 0 &\text{otherwise}\end{cases}) = (-1)^{|t|} \cdot \begin{cases}1 &\text{if } \forall i_{1}\, :\, \iota,\; i_{1} \in t \implies x \in S(i_{1})\\ 0 &\text{otherwise}\end{cases}\).
By \paperlink{premise:6d756c5f697465}{Prerequisite~P\ref{premise:6d756c5f697465}} with \(\mathbb{Z}\) as its type parameter \(\alpha\) and \(P = \forall i_{1}\, :\, \iota,\; i_{1} \in t \implies x \in S\left(i_{1}\right)\) and \(\left(-1\right)^{\left|t\right|}\) in place of its parameter \(a\) and \(b = 1\) and \(c = 0\), we obtain \((-1)^{|t|} \cdot \begin{cases}1 &\text{if } \forall i_{1}\, :\, \iota,\; i_{1} \in t \implies x \in S(i_{1})\\ 0 &\text{otherwise}\end{cases} = \begin{cases}(-1)^{|t|} \cdot 1 &\text{if } \forall i_{1}\, :\, \iota,\; i_{1} \in t \implies x \in S(i_{1})\\ (-1)^{|t|} \cdot 0 &\text{otherwise}\end{cases}\).
Chaining the equalities, we obtain \((-1)^{|t|} \cdot (1 \cdot \begin{cases}1 &\text{if } x \in \bigcap_{x_{1}\in t} S(x_{1})\\ 0 &\text{otherwise}\end{cases}) = \begin{cases}(-1)^{|t|} \cdot 1 &\text{if } \forall i_{1}\, :\, \iota,\; i_{1} \in t \implies x \in S(i_{1})\\ (-1)^{|t|} \cdot 0 &\text{otherwise}\end{cases}\).
By definition, we obtain \((\forall i_{1}\, :\, \iota,\; i_{1} \in t \implies x \in S(i_{1})) = (\forall i_{1}\, :\, \iota,\; i_{1} \in t \implies x \in S(i_{1}))\).
For the required implication, argue as follows.
\begin{proofstep}
Assume \(\forall i_{1}\, :\, \iota,\; i_{1} \in t \implies x \in S\left(i_{1}\right)\).
By \paperlink{premise:6d756c5f6f6e65}{Prerequisite~P\ref{premise:6d756c5f6f6e65}} with \(M = \mathbb{Z}\) and \(\left(-1\right)^{\left|t\right|}\) in place of its parameter \(a\), we obtain \(\left(-1\right)^{\left|t\right|} \cdot 1 = \left(-1\right)^{\left|t\right|}\).
\end{proofstep}
For the required implication, argue as follows.
\begin{proofstep}
Assume \(\neg \left(\forall i_{1}\, :\, \iota,\; i_{1} \in t \implies x \in S\left(i_{1}\right)\right)\).
By \paperlink{premise:4d756c5a65726f436c6173732e6d756c5f7a65726f}{Prerequisite~P\ref{premise:4d756c5a65726f436c6173732e6d756c5f7a65726f}} with \(M_{0} = \mathbb{Z}\) and \(\left(-1\right)^{\left|t\right|}\) in place of its parameter \(a\), we obtain \(\left(-1\right)^{\left|t\right|} \cdot 0 = 0\).
\end{proofstep}
By \paperlink{premise:6974655f636f6e6772}{Prerequisite~P\ref{premise:6974655f636f6e6772}} with \(\mathbb{Z}\) as its type parameter \(\alpha\), we obtain \(\begin{cases}(-1)^{|t|} \cdot 1 &\text{if } \forall i_{1}\, :\, \iota,\; i_{1} \in t \implies x \in S(i_{1})\\ (-1)^{|t|} \cdot 0 &\text{otherwise}\end{cases} = \begin{cases}(-1)^{|t|} &\text{if } \forall i_{1}\, :\, \iota,\; i_{1} \in t \implies x \in S(i_{1})\\ 0 &\text{otherwise}\end{cases}\).
Chaining the equalities, we obtain \((-1)^{|t|} \cdot (1 \cdot \begin{cases}1 &\text{if } x \in \bigcap_{x_{1}\in t} S(x_{1})\\ 0 &\text{otherwise}\end{cases}) = \begin{cases}(-1)^{|t|} &\text{if } \forall i_{1}\, :\, \iota,\; i_{1} \in t \implies x \in S(i_{1})\\ 0 &\text{otherwise}\end{cases}\).
Applying the same function to equal arguments, we obtain \((-1)^{|t|} \cdot (1 \cdot \begin{cases}1 &\text{if } x \in \bigcap_{x_{1}\in t} S(x_{1})\\ 0 &\text{otherwise}\end{cases}) \mathbin{\bullet} f(x) = \begin{cases}(-1)^{|t|} &\text{if } \forall i_{1}\, :\, \iota,\; i_{1} \in t \implies x \in S(i_{1})\\ 0 &\text{otherwise}\end{cases} \mathbin{\bullet} f(x)\).
By \paperlink{premise:6974655f736d756c}{Prerequisite~P\ref{premise:6974655f736d756c}} with \(\beta = \mathbb{Z}\) and \(p = \forall i_{1}\, :\, \iota,\; i_{1} \in t \implies x \in S\left(i_{1}\right)\) and \(f\left(x\right)\) in place of its parameter \(a\) and \(b = \left(-1\right)^{\left|t\right|}\) and \(c = 0\), we obtain \(\begin{cases}(-1)^{|t|} &\text{if } \forall i_{1}\, :\, \iota,\; i_{1} \in t \implies x \in S(i_{1})\\ 0 &\text{otherwise}\end{cases} \mathbin{\bullet} f(x) = \begin{cases}(-1)^{|t|} \mathbin{\bullet} f(x) &\text{if } \forall i_{1}\, :\, \iota,\; i_{1} \in t \implies x \in S(i_{1})\\ 0 \mathbin{\bullet} f(x) &\text{otherwise}\end{cases}\).
Chaining the equalities, we obtain \((-1)^{|t|} \cdot (1 \cdot \begin{cases}1 &\text{if } x \in \bigcap_{x_{1}\in t} S(x_{1})\\ 0 &\text{otherwise}\end{cases}) \mathbin{\bullet} f(x) = \begin{cases}(-1)^{|t|} \mathbin{\bullet} f(x) &\text{if } \forall i_{1}\, :\, \iota,\; i_{1} \in t \implies x \in S(i_{1})\\ 0 \mathbin{\bullet} f(x) &\text{otherwise}\end{cases}\).
By definition, we obtain \((\forall i_{1}\, :\, \iota,\; i_{1} \in t \implies x \in S(i_{1})) = (\forall i_{1}\, :\, \iota,\; i_{1} \in t \implies x \in S(i_{1}))\).
For the required implication, argue as follows.
\begin{proofstep}
Assume \(\forall i_{1}\, :\, \iota,\; i_{1} \in t \implies x \in S\left(i_{1}\right)\).
By definition, we obtain \((-1)^{|t|} \mathbin{\bullet} f(x) = (-1)^{|t|} \mathbin{\bullet} f(x)\).
\end{proofstep}
For the required implication, argue as follows.
\begin{proofstep}
Assume \(\neg \left(\forall i_{1}\, :\, \iota,\; i_{1} \in t \implies x \in S\left(i_{1}\right)\right)\).
By \paperlink{premise:7a65726f5f736d756c}{Prerequisite~P\ref{premise:7a65726f5f736d756c}} with \(M_{0} = \mathbb{Z}\) and \(m = f\left(x\right)\), we obtain \(0 \mathbin{\bullet} f\left(x\right) = 0\).
\end{proofstep}
By \paperlink{premise:6974655f636f6e6772}{Prerequisite~P\ref{premise:6974655f636f6e6772}}, we obtain \(\begin{cases}(-1)^{|t|} \mathbin{\bullet} f(x) &\text{if } \forall i_{1}\, :\, \iota,\; i_{1} \in t \implies x \in S(i_{1})\\ 0 \mathbin{\bullet} f(x) &\text{otherwise}\end{cases} = \begin{cases}(-1)^{|t|} \mathbin{\bullet} f(x) &\text{if } \forall i_{1}\, :\, \iota,\; i_{1} \in t \implies x \in S(i_{1})\\ 0 &\text{otherwise}\end{cases}\).
Chaining the equalities, we obtain \((-1)^{|t|} \cdot (1 \cdot \begin{cases}1 &\text{if } x \in \bigcap_{x_{1}\in t} S(x_{1})\\ 0 &\text{otherwise}\end{cases}) \mathbin{\bullet} f(x) = \begin{cases}(-1)^{|t|} \mathbin{\bullet} f(x) &\text{if } \forall i_{1}\, :\, \iota,\; i_{1} \in t \implies x \in S(i_{1})\\ 0 &\text{otherwise}\end{cases}\).
\end{proofstep}
The substitution uses the following argument.
\begin{proofstep}
Fix \(\mathit{inst}_{2}\) in \(\operatorname{Decidable\,Pred}(x \mapsto \neg (\forall i_{1}\, :\, \iota,\; i_{1} \in t \implies x \in S(i_{1})))\).
For the required implication, argue as follows.
\begin{proofstep}
Fix \(s_{1}\) in \(\mathcal{P}_{\mathrm{fin}}(\alpha)\).
Fix \(s_{2}\) in \(\mathcal{P}_{\mathrm{fin}}(\alpha)\).
Assume \(s_{1} = s_{2}\).
By definition, we obtain \(\{x\in s_{1} \mid \neg (\forall i_{1}\, :\, \iota,\; i_{1} \in t \implies x \in S(i_{1}))\} = \{x\in s_{1} \mid \neg (\forall i_{1}\, :\, \iota,\; i_{1} \in t \implies x \in S(i_{1}))\}\).
Use the hypothesis \(s_{1} = s_{2}\).
Substituting equals, we obtain \(\{x\in s_{1} \mid \neg (\forall i_{1}\, :\, \iota,\; i_{1} \in t \implies x \in S(i_{1}))\} = \{x\in s_{2} \mid \neg (\forall i_{1}\, :\, \iota,\; i_{1} \in t \implies x \in S(i_{1}))\}\).
\end{proofstep}
For the required implication, argue as follows.
\begin{proofstep}
Fix \(a_{1}\) in \(\alpha\).
By \paperlink{premise:696e737453756273696e676c65746f6e446563696461626c65}{Prerequisite~P\ref{premise:696e737453756273696e676c65746f6e446563696461626c65}} with \(p = (x \mapsto \neg (\forall i_{1}\, :\, \iota,\; i_{1} \in t \implies x \in S(i_{1})))(a_{1})\), we obtain \(\operatorname{Decidable}((x \mapsto \neg (\forall i_{1}\, :\, \iota,\; i_{1} \in t \implies x \in S(i_{1})))(a_{1}))\text{ has at most one element}\).
\end{proofstep}
By \paperlink{premise:50692e696e737453756273696e676c65746f6e}{Prerequisite~P\ref{premise:50692e696e737453756273696e676c65746f6e}}, we obtain \((\Pi_{a_{1}\, :\, \alpha}\, (a_{2} \mapsto \operatorname{Decidable}((x \mapsto \neg (\forall i_{1}\, :\, \iota,\; i_{1} \in t \implies x \in S(i_{1})))(a_{2})))(a_{1}))\text{ has at most one element}\).
By \paperlink{premise:53756273696e676c65746f6e2e656c696d}{Prerequisite~P\ref{premise:53756273696e676c65746f6e2e656c696d}} with \(\operatorname{Decidable\,Pred}(x \mapsto \neg (\forall i_{1}\, :\, \iota,\; i_{1} \in t \implies x \in S(i_{1})))\) as its type parameter \(\alpha\) and \(\mathit{inst}_{2}\) in place of its parameter \(a\) and \(b = \mathit{inst}_{3}\), we obtain \(\mathit{inst}_{2} = \mathit{inst}_{3}\).
Substituting equals, we obtain \((\mathit{inst}_{4} \mapsto \forall s_{1}\, :\, \mathcal{P}_{\mathrm{fin}}(\alpha),\; \forall s_{2}\, :\, \mathcal{P}_{\mathrm{fin}}(\alpha),\; s_{1} = s_{2} \implies \{x\in s_{1} \mid \neg (\forall i_{1}\, :\, \iota,\; i_{1} \in t \implies x \in S(i_{1}))\} = \{x\in s_{2} \mid \neg (\forall i_{1}\, :\, \iota,\; i_{1} \in t \implies x \in S(i_{1}))\})(\mathit{inst}_{3})\).
\end{proofstep}
The substitution uses the following argument.
\begin{proofstep}
Fix \(x\) in \(\alpha\).
By \paperlink{premise:436c6173736963616c2e6e6f745f666f72616c6c}{Prerequisite~P\ref{premise:436c6173736963616c2e6e6f745f666f72616c6c}}, we obtain \(\neg (\forall x_{1}\, :\, \iota,\; (x_{2} \mapsto x_{2} \in t \implies x \in S(x_{2}))(x_{1})) \iff (\exists x_{1}\, :\, \iota,\; \neg (x_{2} \mapsto x_{2} \in t \implies x \in S(x_{2}))(x_{1}))\).
By propositional extensionality, we obtain \((\neg (\forall x_{1}\, :\, \iota,\; x_{1} \in t \implies x \in S(x_{1}))) = (\exists x_{1}\, :\, \iota,\; \neg (x_{1} \in t \implies x \in S(x_{1})))\).
For the required implication, argue as follows.
\begin{proofstep}
Fix \(x_{1}\) in \(\iota\).
By \paperlink{premise:436c6173736963616c2e6e6f745f666f72616c6c}{Prerequisite~P\ref{premise:436c6173736963616c2e6e6f745f666f72616c6c}} with \(x_{1} \in t\) as its type parameter \(\alpha\), we obtain \(\neg (\forall x_{2}\, :\, x_{1} \in t,\; (x_{3} \mapsto x \in S(x_{1}))(x_{2})) \iff (\exists x_{2}\, :\, x_{1} \in t,\; \neg (x_{3} \mapsto x \in S(x_{1}))(x_{2}))\).
By propositional extensionality, we obtain \((\neg (x_{1} \in t \implies x \in S(x_{1}))) = (x_{1} \in t \land x \notin S(x_{1}))\).
\end{proofstep}
By function extensionality, we obtain \((x_{1} \mapsto \neg (x_{1} \in t \implies x \in S(x_{1}))) = (x_{1} \mapsto x_{1} \in t \land x \notin S(x_{1}))\).
Applying the same function to equal arguments, we obtain \((\exists x_{1}\, :\, \iota,\; \neg (x_{1} \in t \implies x \in S(x_{1}))) = (\exists x_{1}\, :\, \iota,\; x_{1} \in t \land x \notin S(x_{1}))\).
Chaining the equalities, we obtain \((\neg (\forall x_{1}\, :\, \iota,\; x_{1} \in t \implies x \in S(x_{1}))) = (\exists x_{1}\, :\, \iota,\; x_{1} \in t \land x \notin S(x_{1}))\).
\end{proofstep}
The substitution uses the following argument.
\begin{proofstep}
Fix \(x\) in \(\alpha\).
Assume \(x \in \{x_{1}\in \bigcup_{i_{1}\in s} S(i_{1}) \mid \exists x_{2}\, :\, \iota,\; x_{2} \in t \land x_{1} \notin S(x_{2})\}\).
By definition, we obtain \(0 = 0\).
\end{proofstep}
For the required implication, argue as follows.
\begin{proofstep}
Fix \(f_{1}\) in \(\alpha \to G\).
Fix \(f_{2}\) in \(\alpha \to G\).
Assume \(f_{1} = f_{2}\).
By definition, we obtain \(\sum_{i_{1}\in \bigcap_{i_{1}\in t} S(i_{1})}(f_{1}(i_{1})) = \sum_{i_{1}\in \bigcap_{i_{1}\in t} S(i_{1})}(f_{1}(i_{1}))\).
Use the hypothesis \(f_{1} = f_{2}\).
Substituting equals, we obtain \(\sum_{i_{1}\in \bigcap_{i_{1}\in t} S(i_{1})}(f_{1}(i_{1})) = \sum_{i_{1}\in \bigcap_{i_{1}\in t} S(i_{1})}(f_{2}(i_{1}))\).
\end{proofstep}
For the required implication, argue as follows.
\begin{proofstep}
Fix \(a_{1}\) in \(\alpha\).
The substitution uses the following argument.
\begin{proofstep}
Choose the witness \(i\).
Use the hypothesis \(i \in t\).
This witness proves existence. Hence \(\exists x\, :\, \iota,\; x \in t\).
\end{proofstep}
Assume \(\forall i_{1}\, :\, \iota,\; i_{1} \in t \implies a_{1} \in S\left(i_{1}\right)\).
Choose the witness \(i\).
By \paperlink{premise:46696e7365742e6d656d5f706f776572736574}{Prerequisite~P\ref{premise:46696e7365742e6d656d5f706f776572736574}}, we obtain \(t \in \mathcal P(s) \iff t \subseteq s\).
By propositional extensionality, we obtain \(\left(t \in \mathcal P(s)\right) = \left(t \subseteq s\right)\).
Use the hypothesis \(t \in \mathcal P(s)\).
Use the hypothesis \(i \in t\).
Substituting along the equality, we obtain \(i \in s\).
Rewriting with \(\forall i_{1}\, :\, \iota,\; i_{1} \in t \implies a_{1} \in S\left(i_{1}\right)\) and \(i \in t\) gives \(a_{1} \in S\left(i\right)\).
Combining the two conclusions, we obtain \(i \in s \land a_{1} \in S\left(i\right)\).
This witness proves existence. Hence \(\exists a_{3}\, :\, \iota,\; a_{3} \in s \land a_{1} \in S\left(a_{3}\right)\).
Rewriting with \paperlink{premise:46696e7365742e6d656d5f696e6627}{Prerequisite~P\ref{premise:46696e7365742e6d656d5f696e6627}}, \paperlink{premise:46696e7365742e6d656d5f66696c746572}{Prerequisite~P\ref{premise:46696e7365742e6d656d5f66696c746572}}, \paperlink{premise:46696e7365742e6d656d5f6269556e696f6e}{Prerequisite~P\ref{premise:46696e7365742e6d656d5f6269556e696f6e}}, \paperlink{premise:6966665f616e645f73656c66}{Prerequisite~P\ref{premise:6966665f616e645f73656c66}} and the auxiliary result established above in the preceding result gives \(a_{1} \in \bigcap_{i_{1}\in t} S(i_{1}) \iff a_{1} \in \{x\in \bigcup_{i_{1}\in s} S(i_{1}) \mid \forall i_{1}\, :\, \iota,\; i_{1} \in t \implies x \in S(i_{1})\}\).
\end{proofstep}
By \paperlink{premise:46696e7365742e657874}{Prerequisite~P\ref{premise:46696e7365742e657874}}, we obtain \(\bigcap_{i_{1}\in t} S(i_{1}) = \{x\in \bigcup_{i_{1}\in s} S(i_{1}) \mid \forall i_{1}\, :\, \iota,\; i_{1} \in t \implies x \in S(i_{1})\}\).
By definition, we obtain \((x \mapsto (-1)^{|t|} \mathbin{\bullet} f(x)) = (x \mapsto (-1)^{|t|} \mathbin{\bullet} f(x))\).
Substituting equals, we obtain \(\sum_{x\in \bigcap_{i_{1}\in t} S(i_{1})}((-1)^{|t|} \mathbin{\bullet} f(x)) = \sum_{x\in \{x\in \bigcup_{i_{1}\in s} S(i_{1}) \mid \forall i_{1}\, :\, \iota,\; i_{1} \in t \implies x \in S(i_{1})\}}((-1)^{|t|} \mathbin{\bullet} f(x))\).
Rewriting with \paperlink{premise:46696e7365742e736d756c5f73756d}{Prerequisite~P\ref{premise:46696e7365742e736d756c5f73756d}}, \paperlink{premise:46696e7365742e73756d5f636f6e6772}{Prerequisite~P\ref{premise:46696e7365742e73756d5f636f6e6772}}, \paperlink{premise:46696e7365742e73756d5f697465}{Prerequisite~P\ref{premise:46696e7365742e73756d5f697465}}, \paperlink{premise:66756e657874}{\emph{Function extensionality}} (\paperlink{premise:66756e657874}{Prerequisite~P\ref{premise:66756e657874}}), \paperlink{premise:46696e7365742e73756d5f636f6e73745f7a65726f}{Prerequisite~P\ref{premise:46696e7365742e73756d5f636f6e73745f7a65726f}}, \paperlink{premise:6164645f7a65726f}{Prerequisite~P\ref{premise:6164645f7a65726f}}, the 3 identities established above and the auxiliary result established above in the preceding result gives \((-1)^{|t|} \mathbin{\bullet} \sum_{a_{1}\in \bigcap_{i_{1}\in t} S(i_{1})}(f(a_{1})) = \sum_{x\in \bigcup_{i_{1}\in s} S(i_{1})}((-1)^{|t|} \cdot (1 \cdot \mathbf{1}_{\bigcap_{x_{1}\in t} S(x_{1})}(x)) \mathbin{\bullet} f(x))\).
Rewriting with \paperlink{premise:46696e7365742e73756d5f636f6e6772}{Prerequisite~P\ref{premise:46696e7365742e73756d5f636f6e6772}} and the identity established above in the preceding result gives \((-1)^{|t|} \mathbin{\bullet} \sum_{a_{1}\in \bigcap_{i_{1}\in t} S(i_{1})}(f(a_{1})) = \sum_{x\in \bigcup_{i_{1}\in s} S(i_{1})}((-1)^{|t|} \cdot (\prod_{i_{1}\in s \setminus t}(1) \cdot \prod_{i_{1}\in t}(\mathbf{1}_{S(i_{1})}(x))) \mathbin{\bullet} f(x))\).
\end{proofstep}
Using the existential hypothesis, we obtain \((-1)^{|t|} \mathbin{\bullet} \sum_{a_{1}\in \bigcap_{i\in t} S(i)}(f(a_{1})) = \sum_{x\in \bigcup_{i\in s} S(i)}((-1)^{|t|} \cdot (\prod_{i\in s \setminus t}(1) \cdot \prod_{i\in t}(\mathbf{1}_{S(i)}(x))) \mathbin{\bullet} f(x))\).
Rewriting with \paperlink{premise:6469665f706f73}{Prerequisite~P\ref{premise:6469665f706f73}} (using \(t\ne\varnothing\)) in the preceding result gives \((-1)^{|t|} \mathbin{\bullet} \begin{cases}\sum_{a_{1}\in \bigcap_{i\in t} S(i)}(f(a_{1})) &\text{given } \mathit{ht}_{1}\text{ proving } t\ne\varnothing\\ \sum_{a_{1}\in \bigcup_{i\in s} S(i)}(f(a_{1})) &\text{if } \neg (t\ne\varnothing)\end{cases} = \sum_{x\in \bigcup_{i\in s} S(i)}((-1)^{|t|} \cdot (\prod_{i\in s \setminus t}(1) \cdot \prod_{i\in t}(\mathbf{1}_{S(i)}(x))) \mathbin{\bullet} f(x))\).
\end{proofstep}
Consider the following case.
\begin{proofstep}
Assume \(\neg \left(t\ne\varnothing\right)\).
For the required implication, argue as follows.
\begin{proofstep}
Assume \(\varnothing \in \mathcal P(s)\).
Assume \(\neg \left(\varnothing\ne\varnothing\right)\).
The substitution uses the following argument.
\begin{proofstep}
Fix \(x\) in \(\alpha\).
Assume \(x \in \bigcup_{i\in s} S(i)\).
By \paperlink{premise:706f775f7a65726f}{Prerequisite~P\ref{premise:706f775f7a65726f}} with \(M = \mathbb{Z}\) and \(-1\) in place of its parameter \(a\), we obtain \(\left(-1\right)^{0} = 1\).
Applying the same function to equal arguments, we obtain \(\left(y \mapsto \left(-1\right)^{0} \cdot y\right) = \left(y \mapsto 1 \cdot y\right)\).
By \paperlink{premise:46696e7365742e73646966665f656d707479}{Prerequisite~P\ref{premise:46696e7365742e73646966665f656d707479}}, we obtain \(s \setminus \varnothing = s\).
For the required implication, argue as follows.
\begin{proofstep}
Fix \(x_{1}\) in \(\iota\).
Assume \(x_{1} \in s\).
By definition, we obtain \(1 = 1\).
\end{proofstep}
By \paperlink{premise:46696e7365742e70726f645f636f6e6772}{Prerequisite~P\ref{premise:46696e7365742e70726f645f636f6e6772}} with \(M = \mathbb{Z}\), we obtain \(\prod_{i\in s \setminus \varnothing}\left(1\right) = \prod_{x_{1}\in s}\left(1\right)\).
By \paperlink{premise:46696e7365742e70726f645f636f6e73745f6f6e65}{Prerequisite~P\ref{premise:46696e7365742e70726f645f636f6e73745f6f6e65}} with \(\beta = \mathbb{Z}\), we obtain \(\prod_{x_{1}\in s}\left(1\right) = 1\).
Chaining the equalities, we obtain \(\prod_{i\in s \setminus \varnothing}\left(1\right) = 1\).
Applying the same function to equal arguments, we obtain \(\prod_{i\in s \setminus \varnothing}\left(1\right) \cdot 1 = 1 \cdot 1\).
By \paperlink{premise:6d756c5f6f6e65}{Prerequisite~P\ref{premise:6d756c5f6f6e65}} with \(M = \mathbb{Z}\) and \(1\) in place of its parameter \(a\), we obtain \(1 \cdot 1 = 1\).
Chaining the equalities, we obtain \(\prod_{i\in s \setminus \varnothing}\left(1\right) \cdot 1 = 1\).
By congruence, we obtain \(\left(-1\right)^{0} \cdot \left(\prod_{i\in s \setminus \varnothing}\left(1\right) \cdot 1\right) = 1 \cdot 1\).
By \paperlink{premise:6d756c5f6f6e65}{Prerequisite~P\ref{premise:6d756c5f6f6e65}} with \(M = \mathbb{Z}\) and \(1\) in place of its parameter \(a\), we obtain \(1 \cdot 1 = 1\).
Chaining the equalities, we obtain \(\left(-1\right)^{0} \cdot \left(\prod_{i\in s \setminus \varnothing}\left(1\right) \cdot 1\right) = 1\).
Applying the same function to equal arguments, we obtain \((-1)^{0} \cdot (\prod_{i\in s \setminus \varnothing}(1) \cdot 1) \mathbin{\bullet} f(x) = 1 \mathbin{\bullet} f(x)\).
By \paperlink{premise:6f6e655f736d756c}{Prerequisite~P\ref{premise:6f6e655f736d756c}} with \(M = \mathbb{Z}\) and \(b = f\left(x\right)\), we obtain \(1 \mathbin{\bullet} f\left(x\right) = f\left(x\right)\).
Chaining the equalities, we obtain \((-1)^{0} \cdot (\prod_{i\in s \setminus \varnothing}(1) \cdot 1) \mathbin{\bullet} f(x) = f(x)\).
\end{proofstep}
Rewriting with \paperlink{premise:706f775f7a65726f}{Prerequisite~P\ref{premise:706f775f7a65726f}}, \paperlink{premise:6f6e655f736d756c}{Prerequisite~P\ref{premise:6f6e655f736d756c}}, \paperlink{premise:46696e7365742e73756d5f636f6e6772}{Prerequisite~P\ref{premise:46696e7365742e73756d5f636f6e6772}} and the identity established above gives \((-1)^{0} \mathbin{\bullet} \sum_{a_{2}\in \bigcup_{i\in s} S(i)}(f(a_{2})) = \sum_{x\in \bigcup_{i\in s} S(i)}((-1)^{|\varnothing|} \cdot (\prod_{i\in s \setminus \varnothing}(1) \cdot \prod_{i\in \varnothing}(\mathbf{1}_{S(i)}(x))) \mathbin{\bullet} f(x))\).
\end{proofstep}
By \paperlink{premise:46696e7365742e6e6f745f6e6f6e656d7074795f6966665f65715f656d707479}{Prerequisite~P\ref{premise:46696e7365742e6e6f745f6e6f6e656d7074795f6966665f65715f656d707479}}, we obtain \(\neg \left(t\ne\varnothing\right) \iff t = \varnothing\).
Use the hypothesis \(\neg \left(t\ne\varnothing\right)\).
By the forward implication, we obtain \(t = \varnothing\).
Reversing the equality, we obtain \(\varnothing = t\).
Use the hypothesis \(t \in \mathcal P(s)\).
Use the hypothesis \(\neg \left(t\ne\varnothing\right)\).
Substituting equals, we obtain \((-1)^{|t|} \mathbin{\bullet} \sum_{a_{1}\in \bigcup_{i\in s} S(i)}(f(a_{1})) = \sum_{x\in \bigcup_{i\in s} S(i)}((-1)^{|t|} \cdot (\prod_{i\in s \setminus t}(1) \cdot \prod_{i\in t}(\mathbf{1}_{S(i)}(x))) \mathbin{\bullet} f(x))\).
Rewriting with \paperlink{premise:6469665f6e6567}{Prerequisite~P\ref{premise:6469665f6e6567}} (using \(\neg \left(t\ne\varnothing\right)\)) in the preceding result gives \((-1)^{|t|} \mathbin{\bullet} \begin{cases}\sum_{a_{1}\in \bigcap_{i\in t} S(i)}(f(a_{1})) &\text{given } \mathit{ht}_{1}\text{ proving } t\ne\varnothing\\ \sum_{a_{1}\in \bigcup_{i\in s} S(i)}(f(a_{1})) &\text{if } \neg (t\ne\varnothing)\end{cases} = \sum_{x\in \bigcup_{i\in s} S(i)}((-1)^{|t|} \cdot (\prod_{i\in s \setminus t}(1) \cdot \prod_{i\in t}(\mathbf{1}_{S(i)}(x))) \mathbin{\bullet} f(x))\).
\end{proofstep}
The proposition or its negation holds. Hence \((-1)^{|t|} \mathbin{\bullet} \begin{cases}\sum_{a_{1}\in \bigcap_{i\in t} S(i)}(f(a_{1})) &\text{given } \mathit{ht}\text{ proving } t\ne\varnothing\\ \sum_{a_{1}\in \bigcup_{i\in s} S(i)}(f(a_{1})) &\text{if } \neg (t\ne\varnothing)\end{cases} = \sum_{x\in \bigcup_{i\in s} S(i)}((-1)^{|t|} \cdot (\prod_{i\in s \setminus t}(1) \cdot \prod_{i\in t}(\mathbf{1}_{S(i)}(x))) \mathbin{\bullet} f(x))\).
\end{proofstep}
By \paperlink{premise:46696e7365742e73756d5f636f6e6772}{Prerequisite~P\ref{premise:46696e7365742e73756d5f636f6e6772}} with \(\mathcal{P}_{\mathrm{fin}}(\iota)\) as its type parameter \(\iota\), we obtain \(\sum_{t\in \mathcal P(s)}((-1)^{|t|} \mathbin{\bullet} \begin{cases}\sum_{a\in \bigcap_{i\in t} S(i)}(f(a)) &\text{given } \mathit{ht}\text{ proving } t\ne\varnothing\\ \sum_{a\in \bigcup_{i\in s} S(i)}(f(a)) &\text{if } \neg (t\ne\varnothing)\end{cases}) = \sum_{y\in \mathcal P(s)}(\sum_{x\in \bigcup_{i\in s} S(i)}((-1)^{|y|} \cdot (\prod_{i\in s \setminus y}(1) \cdot \prod_{i\in y}(\mathbf{1}_{S(i)}(x))) \mathbin{\bullet} f(x)))\).
Rewriting with \paperlink{premise:46696e7365742e73756d5f636f6e6772}{Prerequisite~P\ref{premise:46696e7365742e73756d5f636f6e6772}}, \paperlink{premise:46696e7365742e73756d5f636f6d6d}{Prerequisite~P\ref{premise:46696e7365742e73756d5f636f6d6d}} and the identity established above in the preceding result gives \(\sum_{t\in \mathcal P(s)}((-1)^{|t|} \mathbin{\bullet} \begin{cases}\sum_{a\in \bigcap_{i\in t} S(i)}(f(a)) &\text{given } \mathit{ht}\text{ proving } t\ne\varnothing\\ \sum_{a\in \bigcup_{i\in s} S(i)}(f(a)) &\text{if } \neg (t\ne\varnothing)\end{cases}) = \sum_{a\in \bigcup_{i\in s} S(i)}(\prod_{i\in s}(1 - \mathbf{1}_{S(i)}(a)) \mathbin{\bullet} f(a))\).
\par\smallskip\noindent\emph{Equality 4.}\enspace
By definition, we obtain \(\bigcup_{i\in s} S(i) = \bigcup_{i\in s} S(i)\).
For the required implication, argue as follows.
\begin{proofstep}
Fix \(t\) in \(\alpha\).
Assume \(t \in \bigcup_{i\in s} S(i)\).
By definition, we obtain \(s = s\).
For the required implication, argue as follows.
\begin{proofstep}
Fix \(x\) in \(\iota\).
Assume \(x \in s\).
By definition, we obtain \(1 = 1\).
Rewriting with \paperlink{premise:50692e6f6e655f6170706c79}{Prerequisite~P\ref{premise:50692e6f6e655f6170706c79}} in the preceding result gives the same displayed equality.
Rewriting with \paperlink{premise:5365742e696e64696361746f725f6f665f6d656d}{Prerequisite~P\ref{premise:5365742e696e64696361746f725f6f665f6d656d}} (using \(t \in \bigcup_{i\in s} S(i)\)) in the preceding result gives \(1 = \mathbf{1}_{\bigcup_{i\in s} S(i)}(t)\).
For the required implication, argue as follows.
\begin{proofstep}
Assume \(\mathbf{1}_{\bigcup_{i\in s} S(i)}(t) = 1\).
For the required implication, argue as follows.
\begin{proofstep}
Assume \(1 = 1\).
Assume \(\text{the proof of }\left(1 = 1\right)\text{ agrees under transport with the proof of }\left(1 = 1\right)\).
By reflexivity, we obtain \(1 - \mathbf{1}_{S\left(x\right)}(t)\) heterogeneously equal to itself.
\end{proofstep}
Use the hypothesis \(\mathbf{1}_{\bigcup_{i\in s} S(i)}(t) = 1\).
Reversing the equality, we obtain \(1 = \mathbf{1}_{\bigcup_{i\in s} S(i)}(t)\).
By definition, we obtain \(1 = 1\).
Rewriting with \paperlink{premise:50692e6f6e655f6170706c79}{Prerequisite~P\ref{premise:50692e6f6e655f6170706c79}} in the preceding result gives the same displayed equality.
Rewriting with \paperlink{premise:5365742e696e64696361746f725f6f665f6d656d}{Prerequisite~P\ref{premise:5365742e696e64696361746f725f6f665f6d656d}} (using \(t \in \bigcup_{i\in s} S(i)\)) in the preceding result gives \(1 = \mathbf{1}_{\bigcup_{i\in s} S(i)}(t)\).
Substituting equals, we obtain \(\text{the proof of }(1 = \mathbf{1}_{\bigcup_{i\in s} S(i)}(t))\text{ agrees under transport with the proof of }(1 = 1) \implies (1 - \mathbf{1}_{S(x)}(t)\, :\, \mathbb{Z})\text{ is heterogeneously equal to }(\mathbf{1}_{\bigcup_{i\in s} S(i)}(t) - \mathbf{1}_{S(x)}(t)\, :\, \mathbb{Z})\).
\end{proofstep}
By definition, we obtain \(\mathbf{1}_{\bigcup_{i\in s} S(i)}(t) = \mathbf{1}_{\bigcup_{i\in s} S(i)}(t)\).
By definition, we obtain \(1 = 1\).
Rewriting with \paperlink{premise:50692e6f6e655f6170706c79}{Prerequisite~P\ref{premise:50692e6f6e655f6170706c79}} in the preceding result gives the same displayed equality.
Rewriting with \paperlink{premise:5365742e696e64696361746f725f6f665f6d656d}{Prerequisite~P\ref{premise:5365742e696e64696361746f725f6f665f6d656d}} (using \(t \in \bigcup_{i\in s} S(i)\)) in the preceding result gives \(1 = \mathbf{1}_{\bigcup_{i\in s} S(i)}(t)\).
By reflexivity, we obtain the same proof of \(1 = \mathbf{1}_{\bigcup_{i\in s} S(i)}(t)\) on both sides of the heterogeneous equality.
Substituting equals, we obtain \((1 - \mathbf{1}_{S(x)}(t)\, :\, \mathbb{Z})\text{ is heterogeneously equal to }(\mathbf{1}_{\bigcup_{i\in s} S(i)}(t) - \mathbf{1}_{S(x)}(t)\, :\, \mathbb{Z})\).
By \paperlink{premise:65715f6f665f686571}{Prerequisite~P\ref{premise:65715f6f665f686571}} with \(\mathbb{Z}\) as its type parameter \(\alpha\), we obtain \(1 - \mathbf{1}_{S\left(x\right)}(t) = \mathbf{1}_{\bigcup_{i\in s} S(i)}(t) - \mathbf{1}_{S\left(x\right)}(t)\).
\end{proofstep}
By \paperlink{premise:46696e7365742e70726f645f636f6e6772}{Prerequisite~P\ref{premise:46696e7365742e70726f645f636f6e6772}} with \(M = \mathbb{Z}\), we obtain \(\prod_{i\in s}(1 - \mathbf{1}_{S(i)}(t)) = \prod_{i\in s}(\mathbf{1}_{\bigcup_{i_{1}\in s} S(i_{1})}(t) - \mathbf{1}_{S(i)}(t))\).
For the required implication, argue as follows.
\begin{proofstep}
Assume \(\prod_{i\in s}(\mathbf{1}_{\bigcup_{i_{1}\in s} S(i_{1})}(t) - \mathbf{1}_{S(i)}(t)) = \prod_{i\in s}(1 - \mathbf{1}_{S(i)}(t))\).
For the required implication, argue as follows.
\begin{proofstep}
Assume \(\prod_{i\in s}(1 - \mathbf{1}_{S(i)}(t)) = \prod_{i\in s}(1 - \mathbf{1}_{S(i)}(t))\).
Assume \(\text{the proof of }(\prod_{i\in s}(1 - \mathbf{1}_{S(i)}(t)) = \prod_{i\in s}(1 - \mathbf{1}_{S(i)}(t)))\text{ agrees under transport with the proof of }(\prod_{i\in s}(1 - \mathbf{1}_{S(i)}(t)) = \prod_{i\in s}(1 - \mathbf{1}_{S(i)}(t)))\).
By reflexivity, we obtain \(\prod_{i\in s}\left(1 - \mathbf{1}_{S\left(i\right)}(t)\right) \mathbin{\bullet} f\left(t\right)\) heterogeneously equal to itself.
\end{proofstep}
Use the hypothesis \(\prod_{i\in s}(\mathbf{1}_{\bigcup_{i_{1}\in s} S(i_{1})}(t) - \mathbf{1}_{S(i)}(t)) = \prod_{i\in s}(1 - \mathbf{1}_{S(i)}(t))\).
Reversing the equality, we obtain \(\prod_{i\in s}(1 - \mathbf{1}_{S(i)}(t)) = \prod_{i\in s}(\mathbf{1}_{\bigcup_{i_{1}\in s} S(i_{1})}(t) - \mathbf{1}_{S(i)}(t))\).
By definition, we obtain \(s = s\).
For the required implication, argue as follows.
\begin{proofstep}
Fix \(x\) in \(\iota\).
Assume \(x \in s\).
By definition, we obtain \(1 = 1\).
Rewriting with \paperlink{premise:50692e6f6e655f6170706c79}{Prerequisite~P\ref{premise:50692e6f6e655f6170706c79}} in the preceding result gives the same displayed equality.
Rewriting with \paperlink{premise:5365742e696e64696361746f725f6f665f6d656d}{Prerequisite~P\ref{premise:5365742e696e64696361746f725f6f665f6d656d}} (using \(t \in \bigcup_{i\in s} S(i)\)) in the preceding result gives \(1 = \mathbf{1}_{\bigcup_{i\in s} S(i)}(t)\).
For the required implication, argue as follows.
\begin{proofstep}
Assume \(\mathbf{1}_{\bigcup_{i\in s} S(i)}(t) = 1\).
For the required implication, argue as follows.
\begin{proofstep}
Assume \(1 = 1\).
Assume \(\text{the proof of }\left(1 = 1\right)\text{ agrees under transport with the proof of }\left(1 = 1\right)\).
By reflexivity, we obtain \(1 - \mathbf{1}_{S\left(x\right)}(t)\) heterogeneously equal to itself.
\end{proofstep}
Use the hypothesis \(\mathbf{1}_{\bigcup_{i\in s} S(i)}(t) = 1\).
Reversing the equality, we obtain \(1 = \mathbf{1}_{\bigcup_{i\in s} S(i)}(t)\).
By definition, we obtain \(1 = 1\).
Rewriting with \paperlink{premise:50692e6f6e655f6170706c79}{Prerequisite~P\ref{premise:50692e6f6e655f6170706c79}} in the preceding result gives the same displayed equality.
Rewriting with \paperlink{premise:5365742e696e64696361746f725f6f665f6d656d}{Prerequisite~P\ref{premise:5365742e696e64696361746f725f6f665f6d656d}} (using \(t \in \bigcup_{i\in s} S(i)\)) in the preceding result gives \(1 = \mathbf{1}_{\bigcup_{i\in s} S(i)}(t)\).
Substituting equals, we obtain \(\text{the proof of }(1 = \mathbf{1}_{\bigcup_{i\in s} S(i)}(t))\text{ agrees under transport with the proof of }(1 = 1) \implies (1 - \mathbf{1}_{S(x)}(t)\, :\, \mathbb{Z})\text{ is heterogeneously equal to }(\mathbf{1}_{\bigcup_{i\in s} S(i)}(t) - \mathbf{1}_{S(x)}(t)\, :\, \mathbb{Z})\).
\end{proofstep}
By definition, we obtain \(\mathbf{1}_{\bigcup_{i\in s} S(i)}(t) = \mathbf{1}_{\bigcup_{i\in s} S(i)}(t)\).
By definition, we obtain \(1 = 1\).
Rewriting with \paperlink{premise:50692e6f6e655f6170706c79}{Prerequisite~P\ref{premise:50692e6f6e655f6170706c79}} in the preceding result gives the same displayed equality.
Rewriting with \paperlink{premise:5365742e696e64696361746f725f6f665f6d656d}{Prerequisite~P\ref{premise:5365742e696e64696361746f725f6f665f6d656d}} (using \(t \in \bigcup_{i\in s} S(i)\)) in the preceding result gives \(1 = \mathbf{1}_{\bigcup_{i\in s} S(i)}(t)\).
By reflexivity, we obtain the same proof of \(1 = \mathbf{1}_{\bigcup_{i\in s} S(i)}(t)\) on both sides of the heterogeneous equality.
Substituting equals, we obtain \((1 - \mathbf{1}_{S(x)}(t)\, :\, \mathbb{Z})\text{ is heterogeneously equal to }(\mathbf{1}_{\bigcup_{i\in s} S(i)}(t) - \mathbf{1}_{S(x)}(t)\, :\, \mathbb{Z})\).
By \paperlink{premise:65715f6f665f686571}{Prerequisite~P\ref{premise:65715f6f665f686571}} with \(\mathbb{Z}\) as its type parameter \(\alpha\), we obtain \(1 - \mathbf{1}_{S\left(x\right)}(t) = \mathbf{1}_{\bigcup_{i\in s} S(i)}(t) - \mathbf{1}_{S\left(x\right)}(t)\).
\end{proofstep}
By \paperlink{premise:46696e7365742e70726f645f636f6e6772}{Prerequisite~P\ref{premise:46696e7365742e70726f645f636f6e6772}} with \(M = \mathbb{Z}\), we obtain \(\prod_{i\in s}(1 - \mathbf{1}_{S(i)}(t)) = \prod_{i\in s}(\mathbf{1}_{\bigcup_{i_{1}\in s} S(i_{1})}(t) - \mathbf{1}_{S(i)}(t))\).
Substituting equals, we obtain \(\text{the proof of }(\prod_{i\in s}(1 - \mathbf{1}_{S(i)}(t)) = \prod_{i\in s}(\mathbf{1}_{\bigcup_{i_{1}\in s} S(i_{1})}(t) - \mathbf{1}_{S(i)}(t)))\text{ agrees under transport with the proof of }(\prod_{i\in s}(1 - \mathbf{1}_{S(i)}(t)) = \prod_{i\in s}(1 - \mathbf{1}_{S(i)}(t))) \implies (\prod_{i\in s}(1 - \mathbf{1}_{S(i)}(t)) \mathbin{\bullet} f(t)\, :\, G)\text{ is heterogeneously equal to }(\prod_{i\in s}(\mathbf{1}_{\bigcup_{i_{1}\in s} S(i_{1})}(t) - \mathbf{1}_{S(i)}(t)) \mathbin{\bullet} f(t)\, :\, G)\).
\end{proofstep}
By definition, we obtain \(\prod_{i\in s}(\mathbf{1}_{\bigcup_{i_{1}\in s} S(i_{1})}(t) - \mathbf{1}_{S(i)}(t)) = \prod_{i\in s}(\mathbf{1}_{\bigcup_{i_{1}\in s} S(i_{1})}(t) - \mathbf{1}_{S(i)}(t))\).
By definition, we obtain \(s = s\).
For the required implication, argue as follows.
\begin{proofstep}
Fix \(x\) in \(\iota\).
Assume \(x \in s\).
By definition, we obtain \(1 = 1\).
Rewriting with \paperlink{premise:50692e6f6e655f6170706c79}{Prerequisite~P\ref{premise:50692e6f6e655f6170706c79}} in the preceding result gives the same displayed equality.
Rewriting with \paperlink{premise:5365742e696e64696361746f725f6f665f6d656d}{Prerequisite~P\ref{premise:5365742e696e64696361746f725f6f665f6d656d}} (using \(t \in \bigcup_{i\in s} S(i)\)) in the preceding result gives \(1 = \mathbf{1}_{\bigcup_{i\in s} S(i)}(t)\).
For the required implication, argue as follows.
\begin{proofstep}
Assume \(\mathbf{1}_{\bigcup_{i\in s} S(i)}(t) = 1\).
For the required implication, argue as follows.
\begin{proofstep}
Assume \(1 = 1\).
Assume \(\text{the proof of }\left(1 = 1\right)\text{ agrees under transport with the proof of }\left(1 = 1\right)\).
By reflexivity, we obtain \(1 - \mathbf{1}_{S\left(x\right)}(t)\) heterogeneously equal to itself.
\end{proofstep}
Use the hypothesis \(\mathbf{1}_{\bigcup_{i\in s} S(i)}(t) = 1\).
Reversing the equality, we obtain \(1 = \mathbf{1}_{\bigcup_{i\in s} S(i)}(t)\).
By definition, we obtain \(1 = 1\).
Rewriting with \paperlink{premise:50692e6f6e655f6170706c79}{Prerequisite~P\ref{premise:50692e6f6e655f6170706c79}} in the preceding result gives the same displayed equality.
Rewriting with \paperlink{premise:5365742e696e64696361746f725f6f665f6d656d}{Prerequisite~P\ref{premise:5365742e696e64696361746f725f6f665f6d656d}} (using \(t \in \bigcup_{i\in s} S(i)\)) in the preceding result gives \(1 = \mathbf{1}_{\bigcup_{i\in s} S(i)}(t)\).
Substituting equals, we obtain \(\text{the proof of }(1 = \mathbf{1}_{\bigcup_{i\in s} S(i)}(t))\text{ agrees under transport with the proof of }(1 = 1) \implies (1 - \mathbf{1}_{S(x)}(t)\, :\, \mathbb{Z})\text{ is heterogeneously equal to }(\mathbf{1}_{\bigcup_{i\in s} S(i)}(t) - \mathbf{1}_{S(x)}(t)\, :\, \mathbb{Z})\).
\end{proofstep}
By definition, we obtain \(\mathbf{1}_{\bigcup_{i\in s} S(i)}(t) = \mathbf{1}_{\bigcup_{i\in s} S(i)}(t)\).
By definition, we obtain \(1 = 1\).
Rewriting with \paperlink{premise:50692e6f6e655f6170706c79}{Prerequisite~P\ref{premise:50692e6f6e655f6170706c79}} in the preceding result gives the same displayed equality.
Rewriting with \paperlink{premise:5365742e696e64696361746f725f6f665f6d656d}{Prerequisite~P\ref{premise:5365742e696e64696361746f725f6f665f6d656d}} (using \(t \in \bigcup_{i\in s} S(i)\)) in the preceding result gives \(1 = \mathbf{1}_{\bigcup_{i\in s} S(i)}(t)\).
By reflexivity, we obtain the same proof of \(1 = \mathbf{1}_{\bigcup_{i\in s} S(i)}(t)\) on both sides of the heterogeneous equality.
Substituting equals, we obtain \((1 - \mathbf{1}_{S(x)}(t)\, :\, \mathbb{Z})\text{ is heterogeneously equal to }(\mathbf{1}_{\bigcup_{i\in s} S(i)}(t) - \mathbf{1}_{S(x)}(t)\, :\, \mathbb{Z})\).
By \paperlink{premise:65715f6f665f686571}{Prerequisite~P\ref{premise:65715f6f665f686571}} with \(\mathbb{Z}\) as its type parameter \(\alpha\), we obtain \(1 - \mathbf{1}_{S\left(x\right)}(t) = \mathbf{1}_{\bigcup_{i\in s} S(i)}(t) - \mathbf{1}_{S\left(x\right)}(t)\).
\end{proofstep}
By \paperlink{premise:46696e7365742e70726f645f636f6e6772}{Prerequisite~P\ref{premise:46696e7365742e70726f645f636f6e6772}} with \(M = \mathbb{Z}\), we obtain \(\prod_{i\in s}(1 - \mathbf{1}_{S(i)}(t)) = \prod_{i\in s}(\mathbf{1}_{\bigcup_{i_{1}\in s} S(i_{1})}(t) - \mathbf{1}_{S(i)}(t))\).
By reflexivity, we obtain the same proof of \(\prod_{i\in s}(1 - \mathbf{1}_{S(i)}(t)) = \prod_{i\in s}(\mathbf{1}_{\bigcup_{i_{1}\in s} S(i_{1})}(t) - \mathbf{1}_{S(i)}(t))\) on both sides of the heterogeneous equality.
Substituting equals, we obtain \((\prod_{i\in s}(1 - \mathbf{1}_{S(i)}(t)) \mathbin{\bullet} f(t)\, :\, G)\text{ is heterogeneously equal to }(\prod_{i\in s}(\mathbf{1}_{\bigcup_{i_{1}\in s} S(i_{1})}(t) - \mathbf{1}_{S(i)}(t)) \mathbin{\bullet} f(t)\, :\, G)\).
By \paperlink{premise:65715f6f665f686571}{Prerequisite~P\ref{premise:65715f6f665f686571}}, we obtain \(\prod_{i\in s}(1 - \mathbf{1}_{S(i)}(t)) \mathbin{\bullet} f(t) = \prod_{i\in s}(\mathbf{1}_{\bigcup_{i_{1}\in s} S(i_{1})}(t) - \mathbf{1}_{S(i)}(t)) \mathbin{\bullet} f(t)\).
\end{proofstep}
By \paperlink{premise:46696e7365742e73756d5f636f6e6772}{Prerequisite~P\ref{premise:46696e7365742e73756d5f636f6e6772}}, we obtain \(\sum_{a\in \bigcup_{i\in s} S(i)}(\prod_{i\in s}(1 - \mathbf{1}_{S(i)}(a)) \mathbin{\bullet} f(a)) = \sum_{a\in \bigcup_{i\in s} S(i)}(\prod_{i\in s}(\mathbf{1}_{\bigcup_{i_{1}\in s} S(i_{1})}(a) - \mathbf{1}_{S(i)}(a)) \mathbin{\bullet} f(a))\).
\par\smallskip\noindent\emph{Equality 5.}\enspace
By \paperlink{premise:46696e7365742e65715f656d7074795f6f725f6e6f6e656d707479}{Prerequisite~P\ref{premise:46696e7365742e65715f656d7074795f6f725f6e6f6e656d707479}}, we obtain \(s = \varnothing \lor s\ne\varnothing\).
Consider the following case.
\begin{proofstep}
Assume \(s = \varnothing\).
The substitution uses the following argument.
\begin{proofstep}
Fix \(x\) in \(\alpha\).
Assume \(x \in \varnothing\).
By definition, we obtain \(\varnothing = \varnothing\).
For the required implication, argue as follows.
\begin{proofstep}
Fix \(x_{1}\) in \(\iota\).
Assume \(x_{1} \in \varnothing\).
By \paperlink{premise:46696e7365742e636f655f656d707479}{Prerequisite~P\ref{premise:46696e7365742e636f655f656d707479}}, we obtain \(\operatorname{underlying\,set}(\varnothing) = \varnothing\).
Applying the same function to equal arguments, we obtain \(\mathbf{1}_{\varnothing}(x) = \mathbf{1}_{\varnothing}(x)\).
Rewriting with \paperlink{premise:5365742e6d656d5f656d7074795f6966665f66616c7365}{Prerequisite~P\ref{premise:5365742e6d656d5f656d7074795f6966665f66616c7365}} and \paperlink{premise:6e6f745f66616c73655f65715f74727565}{Prerequisite~P\ref{premise:6e6f745f66616c73655f65715f74727565}} gives \(x \notin \varnothing\).
By \paperlink{premise:5365742e696e64696361746f725f6f665f6e6f745f6d656d}{Prerequisite~P\ref{premise:5365742e696e64696361746f725f6f665f6e6f745f6d656d}} with \(M = \mathbb{Z}\) and \(1\) in place of its parameter \(f\), we obtain \(\mathbf{1}_{\varnothing}(x) = 0\).
Chaining the equalities, we obtain \(\mathbf{1}_{\varnothing}(x) = 0\).
Applying the same function to equal arguments, we obtain \(\mathbf{1}_{\varnothing}(x) - \mathbf{1}_{S\left(x_{1}\right)}(x) = 0 - \mathbf{1}_{S\left(x_{1}\right)}(x)\).
By \paperlink{premise:7a65726f5f737562}{Prerequisite~P\ref{premise:7a65726f5f737562}} with \(\mathbb{Z}\) as its type parameter \(G\) and \(\mathbf{1}_{S\left(x_{1}\right)}(x)\) in place of its parameter \(a\), we obtain \(0 - \mathbf{1}_{S\left(x_{1}\right)}(x) = -\mathbf{1}_{S\left(x_{1}\right)}(x)\).
Chaining the equalities, we obtain \(\mathbf{1}_{\varnothing}(x) - \mathbf{1}_{S\left(x_{1}\right)}(x) = -\mathbf{1}_{S\left(x_{1}\right)}(x)\).
\end{proofstep}
By \paperlink{premise:46696e7365742e70726f645f636f6e6772}{Prerequisite~P\ref{premise:46696e7365742e70726f645f636f6e6772}} with \(M = \mathbb{Z}\), we obtain \(\prod_{i\in \varnothing}(\mathbf{1}_{\bigcup_{i_{1}\in \varnothing} S(i_{1})}(x) - \mathbf{1}_{S(i)}(x)) = \prod_{x_{1}\in \varnothing}(-\mathbf{1}_{S(x_{1})}(x))\).
Applying the same function to equal arguments, we obtain \(\prod_{i\in \varnothing}(\mathbf{1}_{\bigcup_{i_{1}\in \varnothing} S(i_{1})}(x) - \mathbf{1}_{S(i)}(x)) \mathbin{\bullet} f(x) = \prod_{x_{1}\in \varnothing}(-\mathbf{1}_{S(x_{1})}(x)) \mathbin{\bullet} f(x)\).
By \paperlink{premise:6f6e655f736d756c}{Prerequisite~P\ref{premise:6f6e655f736d756c}} with \(M = \mathbb{Z}\) and \(b = f\left(x\right)\), we obtain \(1 \mathbin{\bullet} f\left(x\right) = f\left(x\right)\).
Chaining the equalities, we obtain \(\prod_{i\in \varnothing}(\mathbf{1}_{\bigcup_{i_{1}\in \varnothing} S(i_{1})}(x) - \mathbf{1}_{S(i)}(x)) \mathbin{\bullet} f(x) = f(x)\).
\end{proofstep}
Rewriting with \paperlink{premise:46696e7365742e73756d5f636f6e6772}{Prerequisite~P\ref{premise:46696e7365742e73756d5f636f6e6772}} and the identity established above gives \(\sum_{a\in \bigcup_{i\in \varnothing} S(i)}(\prod_{i\in \varnothing}(\mathbf{1}_{\bigcup_{i_{1}\in \varnothing} S(i_{1})}(a) - \mathbf{1}_{S(i)}(a)) \mathbin{\bullet} f(a)) = 0\).
Use the hypothesis \(s = \varnothing\).
Reversing the equality, we obtain \(\varnothing = s\).
Substituting equals, we obtain \(\sum_{a\in \bigcup_{i\in s} S(i)}(\prod_{i\in s}(\mathbf{1}_{\bigcup_{i_{1}\in s} S(i_{1})}(a) - \mathbf{1}_{S(i)}(a)) \mathbin{\bullet} f(a)) = 0\).
\end{proofstep}
Consider the following case.
\begin{proofstep}
Assume \(s\ne\varnothing\).
The substitution uses the following argument.
\begin{proofstep}
Fix \(x\) in \(\alpha\).
Assume \(x \in \bigcup_{i\in s} S(i)\).
By \paperlink{theorem:1}{Lemma~\ref{theorem:1}} with \(x\) in place of its parameter \(a\), we obtain \(\prod_{i\in s}\left(\mathbf{1}_{\bigcup_{i_{1}\in s} S(i_{1})}(x) - \mathbf{1}_{S\left(i\right)}(x)\right) = 0\).
Applying the same function to equal arguments, we obtain \(\prod_{i\in s}(\mathbf{1}_{\bigcup_{i_{1}\in s} S(i_{1})}(x) - \mathbf{1}_{S(i)}(x)) \mathbin{\bullet} f(x) = 0 \mathbin{\bullet} f(x)\).
By \paperlink{premise:7a65726f5f736d756c}{Prerequisite~P\ref{premise:7a65726f5f736d756c}} with \(M_{0} = \mathbb{Z}\) and \(m = f\left(x\right)\), we obtain \(0 \mathbin{\bullet} f\left(x\right) = 0\).
Chaining the equalities, we obtain \(\prod_{i\in s}(\mathbf{1}_{\bigcup_{i_{1}\in s} S(i_{1})}(x) - \mathbf{1}_{S(i)}(x)) \mathbin{\bullet} f(x) = 0\).
\end{proofstep}
Rewriting with \paperlink{premise:46696e7365742e73756d5f636f6e6772}{Prerequisite~P\ref{premise:46696e7365742e73756d5f636f6e6772}}, \paperlink{premise:46696e7365742e73756d5f636f6e73745f7a65726f}{Prerequisite~P\ref{premise:46696e7365742e73756d5f636f6e73745f7a65726f}} and the identity established above gives \(\sum_{a\in \bigcup_{i\in s} S(i)}(\prod_{i\in s}(\mathbf{1}_{\bigcup_{i_{1}\in s} S(i_{1})}(a) - \mathbf{1}_{S(i)}(a)) \mathbin{\bullet} f(a)) = 0\).
\end{proofstep}
The two alternatives exhaust the cases. Hence \(\sum_{a\in \bigcup_{i\in s} S(i)}(\prod_{i\in s}(\mathbf{1}_{\bigcup_{i_{1}\in s} S(i_{1})}(a) - \mathbf{1}_{S(i)}(a)) \mathbin{\bullet} f(a)) = 0\).
\par Combining these equalities gives \(\sum_{a\in \bigcup_{i\in s} S(i)}(f(a)) - \sum_{\varnothing \ne t \subseteq s}((-1)^{|t| + 1} \mathbin{\bullet} \sum_{a\in \bigcap_{i\in t} S(i)}(f(a))) = 0\).
Rewriting with \paperlink{premise:7375625f65715f7a65726f}{Prerequisite~P\ref{premise:7375625f65715f7a65726f}} in the preceding result gives \(\sum_{a\in \bigcup_{i\in s} S(i)}(f(a)) = \sum_{\varnothing \ne t \subseteq s}((-1)^{|t| + 1} \mathbin{\bullet} \sum_{a\in \bigcap_{i\in t} S(i)}(f(a)))\).
\end{proof}

\begin{corollary}[{Inclusion-exclusion principle for the cardinality of a union}]\label{theorem:3}
Let \(\iota\) and \(\alpha\) be types. Assume equality on \(\alpha\) is decidable. Let \(s\) be a finite subset of \(\iota\). Let \(S\) be a family of finite subsets of \(\alpha\) indexed by \(\iota\). Then
\begin{center}
\(\displaystyle |\bigcup_{i\in s} S(i)| = \sum_{\varnothing \ne t \subseteq s}((-1)^{|t| + 1} \cdot |\bigcap_{i\in t} S(i)|).\)
\end{center}
\end{corollary}

\begin{proof}
The substitution uses the following argument.
\begin{proofstep}
Fix \(x\) in \(\alpha\).
Assume \(x \in \bigcup_{i\in s} S(i)\).
By definition, we obtain \(1 = 1\).
\end{proofstep}
The substitution uses the following argument.
\begin{proofstep}
Fix \(x\) in \(\{x_{1}\in \mathcal P(s) \mid x_{1}\ne\varnothing\}\).
Assume \(x \in \operatorname{attach}\left(\{x_{1}\in \mathcal P(s) \mid x_{1}\ne\varnothing\}\right)\).
By definition, we obtain \(\bigcap_{i\in \operatorname{val}(x)} S(i) = \bigcap_{i\in \operatorname{val}(x)} S(i)\).
For the required implication, argue as follows.
\begin{proofstep}
Fix \(x_{1}\) in \(\alpha\).
Assume \(x_{1} \in \bigcap_{i\in \operatorname{val}(x)} S(i)\).
By definition, we obtain \(1 = 1\).
\end{proofstep}
By \paperlink{premise:46696e7365742e73756d5f636f6e6772}{Prerequisite~P\ref{premise:46696e7365742e73756d5f636f6e6772}} with \(M = \mathbb{Z}\), we obtain \(\sum_{a_{1}\in \bigcap_{i\in \operatorname{val}(x)} S(i)}(1) = \sum_{x_{1}\in \bigcap_{i\in \operatorname{val}(x)} S(i)}(1)\).
By \paperlink{premise:46696e7365742e73756d5f636f6e7374}{Prerequisite~P\ref{premise:46696e7365742e73756d5f636f6e7374}} with \(M = \mathbb{Z}\) and \(b = 1\), we obtain \(\sum_{x_{1}\in \bigcap_{i\in \operatorname{val}(x)} S(i)}(1) = |\bigcap_{i\in \operatorname{val}(x)} S(i)| \mathbin{\bullet} 1\).
Chaining the equalities, we obtain \(\sum_{a_{1}\in \bigcap_{i\in \operatorname{val}(x)} S(i)}(1) = |\bigcap_{i\in \operatorname{val}(x)} S(i)| \mathbin{\bullet} 1\).
By \paperlink{premise:6e736d756c5f65715f6d756c}{Prerequisite~P\ref{premise:6e736d756c5f65715f6d756c}} with \(\mathbb{Z}\) as its type parameter \(\alpha\) and \(n = \left|\bigcap_{i\in \operatorname{val}(x)} S(i)\right|\) and \(1\) in place of its parameter \(a\), we obtain \(|\bigcap_{i\in \operatorname{val}(x)} S(i)| \mathbin{\bullet} 1 = \operatorname{cast}_{\mathbb{Z}}(|\bigcap_{i\in \operatorname{val}(x)} S(i)|) \cdot 1\).
By \paperlink{premise:6d756c5f6f6e65}{Prerequisite~P\ref{premise:6d756c5f6f6e65}} with \(M = \mathbb{Z}\) and \(\operatorname{cast}_{\mathbb{Z}}\left(\left|\bigcap_{i\in \operatorname{val}(x)} S(i)\right|\right)\) in place of its parameter \(a\), we obtain \(\operatorname{cast}_{\mathbb{Z}}(|\bigcap_{i\in \operatorname{val}(x)} S(i)|) \cdot 1 = \operatorname{cast}_{\mathbb{Z}}(|\bigcap_{i\in \operatorname{val}(x)} S(i)|)\).
Chaining the equalities, we obtain \(|\bigcap_{i\in \operatorname{val}(x)} S(i)| \mathbin{\bullet} 1 = \operatorname{cast}_{\mathbb{Z}}(|\bigcap_{i\in \operatorname{val}(x)} S(i)|)\).
Chaining the equalities, we obtain \(\sum_{a_{1}\in \bigcap_{i\in \operatorname{val}(x)} S(i)}(1) = \operatorname{cast}_{\mathbb{Z}}(|\bigcap_{i\in \operatorname{val}(x)} S(i)|)\).
Applying the same function to equal arguments, we obtain \((-1)^{|\operatorname{val}(x)| + 1} \mathbin{\bullet} \sum_{a_{1}\in \bigcap_{i\in \operatorname{val}(x)} S(i)}(1) = (-1)^{|\operatorname{val}(x)| + 1} \mathbin{\bullet} \operatorname{cast}_{\mathbb{Z}}(|\bigcap_{i\in \operatorname{val}(x)} S(i)|)\).
\end{proofstep}
By \paperlink{theorem:2}{Lemma~\ref{theorem:2}} with \(G = \mathbb{Z}\) and \(f = 1\), we obtain \(\sum_{a\in \bigcup_{i\in s} S(i)}(1) = \sum_{\varnothing \ne t \subseteq s}((-1)^{|t| + 1} \mathbin{\bullet} \sum_{a\in \bigcap_{i\in t} S(i)}(1))\).
Rewriting with \paperlink{premise:46696e7365742e73756d5f636f6e6772}{Prerequisite~P\ref{premise:46696e7365742e73756d5f636f6e6772}}, \paperlink{premise:46696e7365742e73756d5f636f6e7374}{Prerequisite~P\ref{premise:46696e7365742e73756d5f636f6e7374}}, \paperlink{premise:6e736d756c5f65715f6d756c}{Prerequisite~P\ref{premise:6e736d756c5f65715f6d756c}}, \paperlink{premise:6d756c5f6f6e65}{Prerequisite~P\ref{premise:6d756c5f6f6e65}} and the two identities established above in the preceding result gives \(|\bigcup_{i\in s} S(i)| = \sum_{\varnothing \ne t \subseteq s}((-1)^{|t| + 1} \cdot |\bigcap_{i\in t} S(i)|)\).
\end{proof}

\begin{lemma}[{Inclusion-exclusion principle for the sum of a function over an intersection of complements}]\label{theorem:4}
Let \(\iota\), \(\alpha\), and \(G\) be types. Assume equality on \(\alpha\) is decidable. Equip \(G\) with an abelian group structure. Fix a finite enumeration of \(\alpha\). Let \(s\) be a finite subset of \(\iota\). Let \(S\) be a family of finite subsets of \(\alpha\) indexed by \(\iota\). Let \(f\) be a function of type \(\alpha \to G\). Then
\begin{center}
\(\displaystyle \sum_{a\in \bigcap_{i\in s} S(i)^{\mathrm c}}(f(a)) = \sum_{t\in \mathcal P(s)}((-1)^{|t|} \mathbin{\bullet} \sum_{a\in \bigcap_{i\in t} S(i)}(f(a))).\)
\end{center}
\end{lemma}

\begin{proof}
\par\smallskip\noindent\emph{Equality 1.}\enspace
Rewriting with \paperlink{premise:46696e7365742e73756d5f636f6d706c5f6164645f73756d}{Prerequisite~P\ref{premise:46696e7365742e73756d5f636f6d706c5f6164645f73756d}}, \paperlink{premise:65715f7375625f6966665f6164645f6571}{Prerequisite~P\ref{premise:65715f7375625f6966665f6164645f6571}}, \paperlink{premise:46696e7365742e7375705f65715f6269556e696f6e}{Prerequisite~P\ref{premise:46696e7365742e7375705f65715f6269556e696f6e}} and \paperlink{premise:46696e7365742e636f6d706c5f737570}{Prerequisite~P\ref{premise:46696e7365742e636f6d706c5f737570}} gives \(\sum_{a\in \bigcap_{i\in s} S(i)^{\mathrm c}}(f(a)) = \sum_{a\in \alpha}(f(a)) - \sum_{a\in \bigcup_{i\in s} S(i)}(f(a))\).
\par\smallskip\noindent\emph{Equality 2.}\enspace
The substitution uses the following argument.
\begin{proofstep}
Fix \(x\) in \(\{x_{1}\in \mathcal P(s) \mid x_{1}\ne\varnothing\}\).
Assume \(x \in \operatorname{attach}\left(\{x_{1}\in \mathcal P(s) \mid x_{1}\ne\varnothing\}\right)\).
By \paperlink{premise:706f775f73756363}{Prerequisite~P\ref{premise:706f775f73756363}} with \(M = \mathbb{Z}\) and \(-1\) in place of its parameter \(a\) and \(n = \left|\operatorname{val}(x)\right|\), we obtain \((-1)^{|\operatorname{val}(x)| + 1} = (-1)^{|\operatorname{val}(x)|} \cdot -1\).
By \paperlink{premise:6d756c5f6e6567}{Prerequisite~P\ref{premise:6d756c5f6e6567}} with \(\mathbb{Z}\) as its type parameter \(\alpha\) and \(\left(-1\right)^{\left|\operatorname{val}(x)\right|}\) in place of its parameter \(a\) and \(b = 1\), we obtain \((-1)^{|\operatorname{val}(x)|} \cdot -1 = -((-1)^{|\operatorname{val}(x)|} \cdot 1)\).
By \paperlink{premise:6d756c5f6f6e65}{Prerequisite~P\ref{premise:6d756c5f6f6e65}} with \(M = \mathbb{Z}\) and \(\left(-1\right)^{\left|\operatorname{val}(x)\right|}\) in place of its parameter \(a\), we obtain \((-1)^{|\operatorname{val}(x)|} \cdot 1 = (-1)^{|\operatorname{val}(x)|}\).
Applying the same function to equal arguments, we obtain \(-((-1)^{|\operatorname{val}(x)|} \cdot 1) = -(-1)^{|\operatorname{val}(x)|}\).
Chaining the equalities, we obtain \((-1)^{|\operatorname{val}(x)|} \cdot -1 = -(-1)^{|\operatorname{val}(x)|}\).
Chaining the equalities, we obtain \((-1)^{|\operatorname{val}(x)| + 1} = -(-1)^{|\operatorname{val}(x)|}\).
Applying the same function to equal arguments, we obtain \((y \mapsto (-1)^{|\operatorname{val}(x)| + 1} \mathbin{\bullet} y) = (y \mapsto -(-1)^{|\operatorname{val}(x)|} \mathbin{\bullet} y)\).
By \paperlink{premise:46696e7365742e6d656d5f66696c746572}{Prerequisite~P\ref{premise:46696e7365742e6d656d5f66696c746572}} with \(\mathcal{P}_{\mathrm{fin}}(\iota)\) as its type parameter \(\alpha\), we obtain \(\operatorname{val}(x) \in \{x_{1}\in \mathcal P(s) \mid x_{1}\ne\varnothing\} \iff \operatorname{val}(x) \in \mathcal P(s) \land \operatorname{val}(x)\ne\varnothing\).
By \paperlink{premise:537562747970652e70726f7065727479}{Prerequisite~P\ref{premise:537562747970652e70726f7065727479}} with \(\mathcal{P}_{\mathrm{fin}}(\iota)\) as its type parameter \(\alpha\) and \(\mathit{self} = x\), we obtain \(\operatorname{val}(x) \in \{x_{1}\in \mathcal P(s) \mid x_{1}\ne\varnothing\}\).
By the forward implication, we obtain \(\operatorname{val}(x) \in \mathcal P(s) \land \operatorname{val}(x)\ne\varnothing\).
Taking the second conjunct, we obtain \(\operatorname{val}(x)\ne\varnothing\).
By \paperlink{premise:46696e7365742e696e66275f65715f696e66}{Prerequisite~P\ref{premise:46696e7365742e696e66275f65715f696e66}} with \(\mathcal{P}_{\mathrm{fin}}(\alpha)\) as its type parameter \(\alpha\) and \(S\) in place of its parameter \(f\), we obtain \(\bigcap_{i\in \operatorname{val}(x)} S(i) = \bigcap_{i\in \operatorname{val}(x)} S(i)\).
For the required implication, argue as follows.
\begin{proofstep}
Fix \(x_{1}\) in \(\alpha\).
Assume \(x_{1} \in \bigcap_{i\in \operatorname{val}(x)} S(i)\).
By definition, we obtain \(f\left(x_{1}\right) = f\left(x_{1}\right)\).
\end{proofstep}
By \paperlink{premise:46696e7365742e73756d5f636f6e6772}{Prerequisite~P\ref{premise:46696e7365742e73756d5f636f6e6772}}, we obtain \(\sum_{a_{1}\in \bigcap_{i\in \operatorname{val}(x)} S(i)}(f(a_{1})) = \sum_{x_{1}\in \bigcap_{i\in \operatorname{val}(x)} S(i)}(f(x_{1}))\).
By congruence, we obtain \((-1)^{|\operatorname{val}(x)| + 1} \mathbin{\bullet} \sum_{a_{1}\in \bigcap_{i\in \operatorname{val}(x)} S(i)}(f(a_{1})) = -(-1)^{|\operatorname{val}(x)|} \mathbin{\bullet} \sum_{x_{1}\in \bigcap_{i\in \operatorname{val}(x)} S(i)}(f(x_{1}))\).
By \paperlink{premise:6e65675f736d756c}{Prerequisite~P\ref{premise:6e65675f736d756c}} with \(R = \mathbb{Z}\) and \(r = \left(-1\right)^{\left|\operatorname{val}(x)\right|}\) and \(\sum_{a_{1}\in \bigcap_{i\in \operatorname{val}(x)} S(i)}\left(f\left(a_{1}\right)\right)\) in place of its parameter \(x\), we obtain \(-(-1)^{|\operatorname{val}(x)|} \mathbin{\bullet} \sum_{a_{1}\in \bigcap_{i\in \operatorname{val}(x)} S(i)}(f(a_{1})) = -((-1)^{|\operatorname{val}(x)|} \mathbin{\bullet} \sum_{a_{1}\in \bigcap_{i\in \operatorname{val}(x)} S(i)}(f(a_{1})))\).
Chaining the equalities, we obtain \((-1)^{|\operatorname{val}(x)| + 1} \mathbin{\bullet} \sum_{a_{1}\in \bigcap_{i\in \operatorname{val}(x)} S(i)}(f(a_{1})) = -((-1)^{|\operatorname{val}(x)|} \mathbin{\bullet} \sum_{a_{1}\in \bigcap_{i\in \operatorname{val}(x)} S(i)}(f(a_{1})))\).
\end{proofstep}
The substitution uses the following argument.
\begin{proofstep}
Fix \(\mathit{inst}_{3}\) in \(\operatorname{Decidable\,Pred}\left(x \mapsto \neg \left(x\ne\varnothing\right)\right)\).
For the required implication, argue as follows.
\begin{proofstep}
Fix \(s_{1}\) in \(\mathcal{P}_{\mathrm{fin}}(\mathcal{P}_{\mathrm{fin}}(\iota))\).
Fix \(s_{2}\) in \(\mathcal{P}_{\mathrm{fin}}(\mathcal{P}_{\mathrm{fin}}(\iota))\).
Assume \(s_{1} = s_{2}\).
By definition, we obtain \(\{x\in s_{1} \mid \neg \left(x\ne\varnothing\right)\} = \{x\in s_{1} \mid \neg \left(x\ne\varnothing\right)\}\).
Use the hypothesis \(s_{1} = s_{2}\).
Substituting equals, we obtain \(\{x\in s_{1} \mid \neg \left(x\ne\varnothing\right)\} = \{x\in s_{2} \mid \neg \left(x\ne\varnothing\right)\}\).
\end{proofstep}
For the required implication, argue as follows.
\begin{proofstep}
Fix \(a\) in \(\mathcal{P}_{\mathrm{fin}}(\iota)\).
By \paperlink{premise:696e737453756273696e676c65746f6e446563696461626c65}{Prerequisite~P\ref{premise:696e737453756273696e676c65746f6e446563696461626c65}} with \(p = \left(x \mapsto \neg \left(x\ne\varnothing\right)\right)\left(a\right)\), we obtain \(\operatorname{Decidable}((x \mapsto \neg (x\ne\varnothing))(a))\text{ has at most one element}\).
\end{proofstep}
By \paperlink{premise:50692e696e737453756273696e676c65746f6e}{Prerequisite~P\ref{premise:50692e696e737453756273696e676c65746f6e}} with \(\mathcal{P}_{\mathrm{fin}}(\iota)\) as its type parameter \(\alpha\), we obtain \((\Pi_{a\, :\, \mathcal{P}_{\mathrm{fin}}(\iota)}\, (a_{1} \mapsto \operatorname{Decidable}((x \mapsto \neg (x\ne\varnothing))(a_{1})))(a))\text{ has at most one element}\).
By \paperlink{premise:53756273696e676c65746f6e2e656c696d}{Prerequisite~P\ref{premise:53756273696e676c65746f6e2e656c696d}} with \(\operatorname{Decidable\,Pred}\left(x \mapsto \neg \left(x\ne\varnothing\right)\right)\) as its type parameter \(\alpha\) and \(a = \mathit{inst}_{3}\) and \(b = \mathit{inst}_{4}\), we obtain \(\mathit{inst}_{3} = \mathit{inst}_{4}\).
Substituting equals, we obtain \((\mathit{inst}_{5} \mapsto \forall s_{1}\, :\, \mathcal{P}_{\mathrm{fin}}(\mathcal{P}_{\mathrm{fin}}(\iota)),\; \forall s_{2}\, :\, \mathcal{P}_{\mathrm{fin}}(\mathcal{P}_{\mathrm{fin}}(\iota)),\; s_{1} = s_{2} \implies \{x\in s_{1} \mid \neg (x\ne\varnothing)\} = \{x\in s_{2} \mid \neg (x\ne\varnothing)\})(\mathit{inst}_{4})\).
\end{proofstep}
The substitution uses the following argument.
\begin{proofstep}
Fix \(x\) in \(\mathcal{P}_{\mathrm{fin}}(\iota)\).
By \paperlink{premise:46696e7365742e6e6f745f6e6f6e656d7074795f6966665f65715f656d707479}{Prerequisite~P\ref{premise:46696e7365742e6e6f745f6e6f6e656d7074795f6966665f65715f656d707479}}, we obtain \(\neg \left(x\ne\varnothing\right) \iff x = \varnothing\).
By propositional extensionality, we obtain \(\left(\neg \left(x\ne\varnothing\right)\right) = \left(x = \varnothing\right)\).
\end{proofstep}
The substitution uses the following argument.
\begin{proofstep}
Fix \(x\) in \(\mathcal{P}_{\mathrm{fin}}(\iota)\).
Assume \(x \in \{\varnothing\}\).
By definition, we obtain \((-1)^{|x|} \mathbin{\bullet} \sum_{a_{1}\in \bigcap_{i\in x} S(i)}(f(a_{1})) = (-1)^{|x|} \mathbin{\bullet} \sum_{a_{1}\in \bigcap_{i\in x} S(i)}(f(a_{1}))\).
\end{proofstep}
Rewriting with \paperlink{theorem:2}{Lemma~\ref{theorem:2}}, \paperlink{premise:46696e7365742e73756d5f636f6e6772}{Prerequisite~P\ref{premise:46696e7365742e73756d5f636f6e6772}}, \paperlink{premise:46696e7365742e73756d5f6e65675f64697374726962}{Prerequisite~P\ref{premise:46696e7365742e73756d5f6e65675f64697374726962}}, \paperlink{premise:7375625f65715f6164645f6e6567}{Prerequisite~P\ref{premise:7375625f65715f6164645f6e6567}}, \paperlink{premise:6e65675f6e6567}{Prerequisite~P\ref{premise:6e65675f6e6567}}, \paperlink{premise:66756e657874}{\emph{Function extensionality}} (\paperlink{premise:66756e657874}{Prerequisite~P\ref{premise:66756e657874}}), \paperlink{premise:46696e7365742e66696c7465725f657127}{Prerequisite~P\ref{premise:46696e7365742e66696c7465725f657127}}, \paperlink{premise:6974655f636f6e645f65715f74727565}{Prerequisite~P\ref{premise:6974655f636f6e645f65715f74727565}}, \paperlink{premise:46696e7365742e6d656d5f706f776572736574}{Prerequisite~P\ref{premise:46696e7365742e6d656d5f706f776572736574}}, \paperlink{premise:46696e7365742e656d7074795f737562736574}{Prerequisite~P\ref{premise:46696e7365742e656d7074795f737562736574}}, \paperlink{premise:46696e7365742e73756d5f73696e676c65746f6e}{Prerequisite~P\ref{premise:46696e7365742e73756d5f73696e676c65746f6e}}, \paperlink{premise:706f775f7a65726f}{Prerequisite~P\ref{premise:706f775f7a65726f}}, \paperlink{premise:6f6e655f736d756c}{Prerequisite~P\ref{premise:6f6e655f736d756c}}, \paperlink{premise:46696e7365742e73756d5f617474616368}{Prerequisite~P\ref{premise:46696e7365742e73756d5f617474616368}}, the 3 identities established above and the auxiliary result established above gives \(\sum_{a\in \alpha}(f(a)) - \sum_{a\in \bigcup_{i\in s} S(i)}(f(a)) = \sum_{t\in \{x\in \mathcal P(s) \mid \neg (x\ne\varnothing)\}}((-1)^{|t|} \mathbin{\bullet} \sum_{a\in \bigcap_{i\in t} S(i)}(f(a))) + \sum_{\varnothing \ne x \subseteq s}((-1)^{|x|} \mathbin{\bullet} \sum_{a\in \bigcap_{i\in x} S(i)}(f(a)))\).
\par\smallskip\noindent\emph{Equality 3.}\enspace
By \paperlink{premise:46696e7365742e73756d5f66696c7465725f6e6f745f6164645f73756d5f66696c746572}{Prerequisite~P\ref{premise:46696e7365742e73756d5f66696c7465725f6e6f745f6164645f73756d5f66696c746572}} with \(\mathcal{P}_{\mathrm{fin}}(\iota)\) as its type parameter \(\iota\) and \(\mathcal P(s)\) in place of its parameter \(s\) and \(p = \operatorname{Nonempty}\) and \(t \mapsto (-1)^{|t|} \mathbin{\bullet} \sum_{a\in \bigcap_{i\in t} S(i)}(f(a))\) in place of its parameter \(f\), we obtain \(\sum_{x\in \{x\in \mathcal P(s) \mid \neg (x\ne\varnothing)\}}((-1)^{|x|} \mathbin{\bullet} \sum_{a\in \bigcap_{i\in x} S(i)}(f(a))) + \sum_{\varnothing \ne x \subseteq s}((-1)^{|x|} \mathbin{\bullet} \sum_{a\in \bigcap_{i\in x} S(i)}(f(a))) = \sum_{x\in \mathcal P(s)}((-1)^{|x|} \mathbin{\bullet} \sum_{a\in \bigcap_{i\in x} S(i)}(f(a)))\).
\par Combining these equalities gives \(\sum_{a\in \bigcap_{i\in s} S(i)^{\mathrm c}}(f(a)) = \sum_{t\in \mathcal P(s)}((-1)^{|t|} \mathbin{\bullet} \sum_{a\in \bigcap_{i\in t} S(i)}(f(a)))\).
\end{proof}

\begin{corollary}[{Inclusion-exclusion principle for the cardinality of an intersection of complements}]\label{theorem:5}
Let \(\iota\) and \(\alpha\) be types. Assume equality on \(\alpha\) is decidable. Fix a finite enumeration of \(\alpha\). Let \(s\) be a finite subset of \(\iota\). Let \(S\) be a family of finite subsets of \(\alpha\) indexed by \(\iota\). Then
\begin{center}
\(\displaystyle |\bigcap_{i\in s} S(i)^{\mathrm c}| = \sum_{t\in \mathcal P(s)}((-1)^{|t|} \cdot |\bigcap_{i\in t} S(i)|).\)
\end{center}
\end{corollary}

\begin{proof}
The substitution uses the following argument.
\begin{proofstep}
Fix \(x\) in \(\alpha\).
Assume \(x \in \bigcap_{i\in s} S\left(i\right)^{\mathrm c}\).
By definition, we obtain \(1 = 1\).
\end{proofstep}
The substitution uses the following argument.
\begin{proofstep}
Fix \(x\) in \(\mathcal{P}_{\mathrm{fin}}(\iota)\).
Assume \(x \in \mathcal P(s)\).
By definition, we obtain \(\bigcap_{i\in x} S(i) = \bigcap_{i\in x} S(i)\).
For the required implication, argue as follows.
\begin{proofstep}
Fix \(x_{1}\) in \(\alpha\).
Assume \(x_{1} \in \bigcap_{i\in x} S(i)\).
By definition, we obtain \(1 = 1\).
\end{proofstep}
By \paperlink{premise:46696e7365742e73756d5f636f6e6772}{Prerequisite~P\ref{premise:46696e7365742e73756d5f636f6e6772}} with \(M = \mathbb{Z}\), we obtain \(\sum_{a_{1}\in \bigcap_{i\in x} S(i)}\left(1\right) = \sum_{x_{1}\in \bigcap_{i\in x} S(i)}\left(1\right)\).
By \paperlink{premise:46696e7365742e73756d5f636f6e7374}{Prerequisite~P\ref{premise:46696e7365742e73756d5f636f6e7374}} with \(M = \mathbb{Z}\) and \(b = 1\), we obtain \(\sum_{x_{1}\in \bigcap_{i\in x} S(i)}\left(1\right) = \left|\bigcap_{i\in x} S(i)\right| \mathbin{\bullet} 1\).
Chaining the equalities, we obtain \(\sum_{a_{1}\in \bigcap_{i\in x} S(i)}\left(1\right) = \left|\bigcap_{i\in x} S(i)\right| \mathbin{\bullet} 1\).
By \paperlink{premise:6e736d756c5f65715f6d756c}{Prerequisite~P\ref{premise:6e736d756c5f65715f6d756c}} with \(\mathbb{Z}\) as its type parameter \(\alpha\) and \(n = \left|\bigcap_{i\in x} S(i)\right|\) and \(1\) in place of its parameter \(a\), we obtain \(|\bigcap_{i\in x} S(i)| \mathbin{\bullet} 1 = \operatorname{cast}_{\mathbb{Z}}(|\bigcap_{i\in x} S(i)|) \cdot 1\).
By \paperlink{premise:6d756c5f6f6e65}{Prerequisite~P\ref{premise:6d756c5f6f6e65}} with \(M = \mathbb{Z}\) and \(\operatorname{cast}_{\mathbb{Z}}\left(\left|\bigcap_{i\in x} S(i)\right|\right)\) in place of its parameter \(a\), we obtain \(\operatorname{cast}_{\mathbb{Z}}(|\bigcap_{i\in x} S(i)|) \cdot 1 = \operatorname{cast}_{\mathbb{Z}}(|\bigcap_{i\in x} S(i)|)\).
Chaining the equalities, we obtain \(|\bigcap_{i\in x} S(i)| \mathbin{\bullet} 1 = \operatorname{cast}_{\mathbb{Z}}(|\bigcap_{i\in x} S(i)|)\).
Chaining the equalities, we obtain \(\sum_{a_{1}\in \bigcap_{i\in x} S(i)}(1) = \operatorname{cast}_{\mathbb{Z}}(|\bigcap_{i\in x} S(i)|)\).
Applying the same function to equal arguments, we obtain \((-1)^{|x|} \mathbin{\bullet} \sum_{a_{1}\in \bigcap_{i\in x} S(i)}(1) = (-1)^{|x|} \mathbin{\bullet} \operatorname{cast}_{\mathbb{Z}}(|\bigcap_{i\in x} S(i)|)\).
\end{proofstep}
By \paperlink{theorem:4}{Lemma~\ref{theorem:4}} with \(G = \mathbb{Z}\) and \(f = 1\), we obtain \(\sum_{a\in \bigcap_{i\in s} S(i)^{\mathrm c}}(1) = \sum_{t\in \mathcal P(s)}((-1)^{|t|} \mathbin{\bullet} \sum_{a\in \bigcap_{i\in t} S(i)}(1))\).
Rewriting with \paperlink{premise:46696e7365742e73756d5f636f6e6772}{Prerequisite~P\ref{premise:46696e7365742e73756d5f636f6e6772}}, \paperlink{premise:46696e7365742e73756d5f636f6e7374}{Prerequisite~P\ref{premise:46696e7365742e73756d5f636f6e7374}}, \paperlink{premise:6e736d756c5f65715f6d756c}{Prerequisite~P\ref{premise:6e736d756c5f65715f6d756c}}, \paperlink{premise:6d756c5f6f6e65}{Prerequisite~P\ref{premise:6d756c5f6f6e65}} and the two identities established above in the preceding result gives \(|\bigcap_{i\in s} S(i)^{\mathrm c}| = \sum_{t\in \mathcal P(s)}((-1)^{|t|} \cdot |\bigcap_{i\in t} S(i)|)\).
\end{proof}

\end{document}
